\documentclass[11pt,a4paper]{article}
\usepackage[margin=22mm]{geometry}
\usepackage{fontspec}
\setmainfont{Latin Modern Roman}
\setsansfont{Latin Modern Sans}
\setmonofont[Scale=0.85]{DejaVu Sans Mono}
\usepackage{babel}
\babelprovide[main,import]{czech}
\usepackage{amsmath,amssymb,amsthm,mathtools}
\renewcommand{\proofname}{Důkaz}
\usepackage{booktabs,longtable,array}
\usepackage{enumitem}
\usepackage{xcolor}
\usepackage{hyperref}
\usepackage{xurl}
\usepackage{mathrsfs}
\usepackage{fancyhdr}
\hypersetup{colorlinks=true,linkcolor=blue!45!black,urlcolor=blue!45!black,pdftitle={TWIST-J: vázané větve a skutečné rozptylové prahy},pdfauthor={Analytická pracovní poznámka}}
\setlength{\parindent}{0pt}
\setlength{\parskip}{5pt}
\setlength{\emergencystretch}{3em}
\setlist{nosep,leftmargin=1.6em}
\pagestyle{fancy}
\fancyhf{}
\fancyhead[L]{\small\sffamily TWIST-J / vázané větve a rozptylové prahy}
\fancyhead[R]{\small\sffamily NON-CANONICAL / L1}
\fancyfoot[C]{\small\thepage}
\setlength{\headheight}{15pt}
\newcommand{\norm}[1]{\left\lVert#1\right\rVert}
\newcommand{\abs}[1]{\left\lvert#1\right\rvert}
\newcommand{\ind}{\mathbf1}
\newcommand{\dd}{\,\mathrm d}
\newcommand{\Rea}{\operatorname{Re}}
\newcommand{\Ran}{\operatorname{Ran}}
\newcommand{\Ccal}{\mathcal C}
\newcommand{\Tcal}{\mathcal T}
\newcommand{\eps}{\varepsilon}
\newtheorem{theorem}{Tvrzení}[section]
\newtheorem{lemma}[theorem]{Pomocné tvrzení}
\theoremstyle{remark}\newtheorem{remark}[theorem]{Poznámka}
\numberwithin{equation}{section}
\providecommand{\tightlist}{\setlength{\itemsep}{0pt}\setlength{\parskip}{0pt}}

\begin{document}

\thispagestyle{empty}
{\Large\sffamily\bfseries Skutečné rozptylové prahy TWIST-J\par}
\vspace{3mm}
{\large Vázané větve na celém hybnostním kruhu\par}
\vspace{4mm}
Analytické pokračování, 29. října 2026\par
\textbf{Samostatný návrh důkazu. NON-CANONICAL / L1.}

\section*{Rozsah této poznámky}

Tato práce navazuje na výslovnou konečnou časovou mez pro původní normalizovaný vstup z PR \#1815 a pro úplný rozptylový operátor z PR \#1817. Zkoumá skutečné fyzikální překážky, které hrubý obecný diskriminant pouze obklopoval. Cílem je připravit podstatně silnější časový odhad.

Podařilo se analyticky určit všechny objektové póly viditelné v sudém kontaktu \(A\), tedy všechny bodové větve potřebné v úplném \(F_2\). V obecném vlákně jsou právě dva nebo tři. Každá větev je určena jedinou monotónní fázovou rovnicí. Prahové konce a nulové kontaktní váhy v kompaktních bodech jsou započteny zvlášť.

Hlavní kladná větev nemůže v původním energetickém okně vytvořit stacionární spojení s pulzem. Obě nízké větve mají doložený obor fáze
\[
\frac12<\omega_\ell<\frac23
\]
a velikost rychlosti menší než \(1/2\). Také ony jsou v daném okně nestacionární. Pro poslední zbývající větev jsou doloženy právě dvě stacionární
konfigurace, se společnou křivostí větší než \(1/4\). Obě mají
kladné kontaktní váhy a leží přísně uvnitř objektové mezery.
Jejich kruhová vzdálenost od již klasifikovaných běžných
kontinuálních stacionárních křivek je větší než \(1/300\).

V místech setkání objektového pólu a kontinuální hrany je doložen společný místní holomorfní příspěvek. Pól a kontinuum se zachovávají v jedné správně pokračované diferenciální formě a na jednom vnějším obvodu. Neodstraňuje se žádný fyzikální člen.

\textbf{Tato poznámka sama nevyhlašuje nový menší konečný čas.} K němu ještě potřebujeme společné číselné velikosti příslušných okolí, úplné místní odhady zbývajících prahů a následný Fourierův účet. Dosavadní čas v samostatném interně přezkoumaném návrhu
\[
N_-=10^{16},\qquad N_+=2^{81\,000\,000\,000\,000}
\]
a jeho úplná časová chyba menší než \(2{,}310002\cdot10^{-7}\) se tím nemění.

\vfill
\small
Původní ostrý vstup, pevný počátek času, absolutní normalizace, celý vybraný jev, všech 26 štítků a korelace s pamětí a referencí jsou zachovány. Výsledek prošel vnitřním analytickým přezkumem v této relaci. Není externím přijetím, součástí původního auditu 130 podmínek ani provedeným fyzickým pokusem. Práce zůstává oddělena od rezervace \#1816. Veřejný zápis ani nový vědecký výpočet nebyly provedeny.
\newpage
\normalsize

\section{Původní operátor a přesný význam stacionárního bodu}

Přejímáme původní úplné kontaktní vzorce z [S1,S2], původní paket a chybový účet z [S3]. Jeho celková hybnost \(K\) má míru \(\mathrm dK/(2\pi)\) na celém kruhu. Vybrané malé výstupní okno \(I_K\) tuto zdrojovou míru neomezuje.

Používáme \(g=2\pi/5\), \(\beta=4\pi/5\) a
\[
e=E-\beta\in I_e=
\left[\frac{4\pi}{5}-\frac1{1000},\frac{4\pi}{5}+\frac1{1000}\right],
\qquad k=K-P.
\]
Pulz má fáze \(\sigma\varepsilon(k)\), \(\sigma=\pm1\), kde
\[
\cos\varepsilon(k)=\frac{1-4\cos k}{5},\qquad
0\le\varepsilon\le a_p=\arccos(-3/5),
\]
\[
\varepsilon'(k)=-\frac{2\sin(k/2)}{\sqrt{3-2\cos k}},\qquad
\varepsilon''(k)=-\frac{\cos(k/2)}{(3-2\cos k)^{3/2}}.
\]
Pro objektovou vázanou fázi \(\omega_j(P)\) je stacionární spojení určeno současně rovnicemi
\[
e=\omega_j(P)+\sigma\varepsilon(K-P),\qquad
\omega_j'(P)=\sigma\varepsilon'(K-P).
\]
Rozhodující křivost při pevném \(K\) je
\[
\partial_P^2[\omega_j(P)+\sigma\varepsilon(K-P)]
=\omega_j''(P)+\sigma\varepsilon''(K-P).
\]
Jednoduchost objektového vlastního stavu sama tuto společnou křivost neomezuje. Právě obě rovnice a toto rozlišení se v následujícím důkazu zachovávají.



\subsection{Zachovaný vstup, jev a úplný přenos}

Pro přehled uvádíme i pevnou normalizaci, kterou nové místní
rozbory nemění. V původním značení [S1--S3] je
\[
h=10^{-3},\quad E_*=8\pi/5,\quad I=[E_*-h,E_*+h],
\]
\[
\psi_{\rm obj}
=[2h\rho_{\rm in}(H_1)]^{-1/2}\,
\mathbf1_I(H_1)\mathbf1_{[-h,h]}(P)a_0,\qquad
a_0=\delta_{R=0,r=0}\,iB,\qquad
\|\psi_{\rm obj}\|=1.
\]
Token je přesně v \(y=0\). Původní celková hybnost je tedy
rovnoměrná na celém kruhu. Energetická míra normalizovaného
vstupu je přesně \(\mathbf1_I(E)\,\mathrm dE/(2h)\).
Původní vybraný výstupní projektor zůstává
\[
\begin{aligned}
Q={}&\mathbf1_{I_K}(K)\mathbf1_{|P|\le1/100}
\mathbf1_I\bigl(\Omega_b(P)+\beta+\varepsilon(k)\bigr)\Pi_b(P)\\
&\otimes P_+(k)\mathbf1_{\varepsilon'(k)>0}
\otimes I_{\rm label}\otimes|d_1\rangle\langle d_1|
\otimes I_{\rm ref}.
\end{aligned}
\]
V něm \(I_{\rm label}\) zahrnuje všech 26 štítků. Paměť se
smí vázat s libovolnou referencí. Ani \(Q\), ani zdroj se v této
poznámce nepodmiňuje na nově vybranou oblast hybnosti.

Úplné kontaktní funkce a přenosový sloupec jsou
\[
F_1(E)=\int_{-\pi}^{\pi}
f_B(P,E-\beta)\,\frac{\mathrm dP}{2\pi},
\]
\[
F_2(K,E)=\sum_{\sigma=\pm1}\int_{-\pi}^{\pi}
h_\sigma(K-P)\,
f_A\bigl(P,E-\beta-\sigma\varepsilon(K-P)\bigr)
\frac{\mathrm dP}{2\pi},
\qquad h_\sigma(k)=\frac{1-\sigma\varepsilon'(k)}2,
\]
\[
T(K,E)=\frac{1}{4+F_1F_2}
\begin{pmatrix}-2F_2&-4\\4&-2F_1\end{pmatrix},
\qquad
V(K,E)=T(K,E)e_1=\frac{(-2F_2,4)^T}{4+F_1F_2}.
\]
Všechny funkce se vyhodnocují z původní vnější strany
\(E\mapsto E+i0\). V \(F_2\) se ponechává celé kontinuum,
každý kontaktně viditelný objektový pól i obě pulzní větve.
Právě regularita celého \(V\), se správně započtenými ostrými
konci \(I\), určuje následný časový odhad.

Označení \(\omega_c(v)\) v soupisu a \(\omega_A(p)\) v části
o křivosti popisují stejnou hlavní kladnou větev:
\(\omega_A(p)=\omega_c(\cos p)\).
Symbol \(\beta(v)\) níže je koncová fáze objektové mezery;
konstantní \(\beta=4\pi/5\) nadále označuje původní fázový posun.
Rovnice označené PT jsou převzaté z oddílu 14 předchozího
časového návrhu [S4]. Nová tvrzení zde mají značky
BC, M, L, BS, CS a PE.

\section{Úplný soupis objektových pólů viditelných v kontaktu A}
Zde se klasifikují všechny sudé kontaktně viditelné póly, tedy všechny bodové větve potřebné v \(F_2\). U kompaktních bodů to není tvrzení o celém bodovém spektru objektu. Lichý kontakt \(B\) je zkontrolován odděleně na konci oddílu.

\hypertarget{dvux11b-skuteux10dnuxe9-mezery-a-fyzikuxe1lnuxed-kontaktnuxed-rovnice}{%
\subsection{Dvě skutečné mezery a fyzikální kontaktní
rovnice}\label{dvux11b-skuteux10dnuxe9-mezery-a-fyzikuxe1lnuxed-kontaktnuxed-rovnice}}

Položme

\[
g=\frac{2\pi}{5},\qquad v=\cos P,\qquad
b=\arccos\frac35=2\arctan\frac12,\qquad a=\pi-b.
\tag{BC1}
\]

Pro \(-1<v<1\) definujme

\[
\alpha(v)=\arccos\frac{v+4}{5},\qquad
\beta(v)=\arccos\frac{v-4}{5}.
\tag{BC2}
\]

Přísné fyzikální mezery \(|t|<4\), \(t=v-5\cos\omega\), jsou

\[
I_+(v)=(\alpha(v),\beta(v)),\qquad
I_-(v)=(2\pi-\beta(v),2\pi-\alpha(v)).
\tag{BC3}
\]

Platí \(0<\alpha<b<g<a<\beta<\pi\). V těchto mezerách

\[
\eta(v,\omega)=\operatorname{sgn}(\sin\omega)
\arccos\frac{v-5\cos\omega}{4},\qquad q=e^{-i\eta},
\tag{BC4}
\]

\[
\partial_\omega\eta
=-\frac{5|\sin\omega|}{\sqrt{16-t^2}}<0.
\tag{BC5}
\]

Úplný sudý jmenovatel, zde označený \(D_o\), je

\[
D_o=xC+qB,
\quad \lambda=e^{-i\omega},\quad x=e^{-i(\omega-g)},
\quad C=x^2-\frac25vx-\frac35,
\quad B=\frac35x^2+\frac25vx-1.
\tag{BC6}
\]

Pro \(-1<v<1\) nemůže na jednotkovém kruhu současně nastat \(B=C=0\).
Proto je fyzikální pólem právě řešení

\[
q=q_{\rm pole}
=e^{-i\delta}\frac{\cos\delta-v-4i\sin\delta}
                         {\cos\delta-v+4i\sin\delta},
\qquad \delta=\omega-g.
\tag{BC7}
\]

V původním důkazu je pro \(\sin P\ne0\) dokázáno, že každý fyzikální
sudý vlastní stav v mezeře je jednoduchý a A-viditelný. Nuly \(D_o\) na
této vrstvě se tedy uvnitř obou mezer nekrátí s čitatelem a jsou právě
hledanými póly.

\hypertarget{monotuxf3nnuxed-fuxe1ze-nahrazuje-nefyzikuxe1lnuxed-algebraickuxe9-koux159eny}{%
\subsection{Monotónní fáze nahrazuje nefyzikální algebraické
kořeny}\label{monotuxf3nnuxed-fuxe1ze-nahrazuje-nefyzikuxe1lnuxed-algebraickuxe9-koux159eny}}

Pro pevné \(-1<v<1\) zvolme spojitý rozvinutý argument

\[
\theta_v(\delta)
=\arg(\cos\delta-v+4i\sin\delta),\qquad \theta_v(0)=0.
\tag{BC8}
\]

Elipsa ve výrazu uvnitř argumentu nemá nulu a jednou obíhá počátek.
Tudíž \(\theta_v(\delta+2\pi)=\theta_v(\delta)+2\pi\),
\(\theta_v(-\delta)=-\theta_v(\delta)\) a \(\theta_v(\pi)=\pi\). Přímé
derivování dává

\[
\partial_\delta\theta_v
=\frac{4(1-v\cos\delta)}
       {(\cos\delta-v)^2+16\sin^2\delta}>0,
\qquad
\partial_v\theta_v
=\frac{4\sin\delta}
       {(\cos\delta-v)^2+16\sin^2\delta}.
\tag{BC9}
\]

Definujme

\[
\varphi_v(\delta)=\delta+2\theta_v(\delta),\qquad
F(v,\omega)=\varphi_v(\omega-g)-\eta(v,\omega).
\tag{BC10}
\]

Pak (BC7) je ekvivalentní přesně

\[
\boxed{F(v,\omega)=2\pi m,\qquad m\in\mathbb Z,}
\tag{BC11}
\]

a na každé z obou přísných mezer

\[
\boxed{
F_\omega
=1+\frac{8(1-v\cos\delta)}
          {(\cos\delta-v)^2+16\sin^2\delta}
 +\frac{5|\sin\omega|}{\sqrt{16-t^2}}>1.
}
\tag{BC12}
\]

Každé povolené celé \(m\) tedy dává nejvýše jeden kořen. Rovnice (BC11),
s níže určeným \(m\) a intervalem, je přesnou jednoznačnou parametrizací
kořene. Nevyžaduje hledání všech šesti algebraických kandidátů ani
dodatečnou numerickou volbu vrstvy.

\hypertarget{hraniux10dnuxed-fuxe1ze-a-uxfaplnuxfd-poux10det-koux159enux16f-v-kladnuxe9-mezeux159e}{%
\subsection{Hraniční fáze a úplný počet kořenů v kladné
mezeře}\label{hraniux10dnuxed-fuxe1ze-a-uxfaplnuxfd-poux10det-koux159enux16f-v-kladnuxe9-mezeux159e}}

Označme jednostranné hodnoty

\[
F_L(v)=\varphi_v(\alpha(v)-g)-\pi,\qquad
F_R(v)=\varphi_v(\beta(v)-g).
\tag{BC13}
\]

Z \(-\pi<\alpha-g<0<\beta-g<\pi\) a z rozsahu argumentu plyne

\[
-4\pi<F_L<-\pi,\qquad 0<F_R<3\pi.
\tag{BC14}
\]

Ve skutečnosti platí \(F_R<2\pi\). Kdyby nastalo \(F_R=2\pi\), musel by
jmenovatel nulovat na horní hraně, kde \(q=1\). Tam

\[
D_o=(x-1)\left[x^2+\left(\frac85-\frac25v\right)x+1\right].
\tag{BC15}
\]

První faktor by vyžadoval \(\beta=g\), což je vyloučeno. Druhý by dal
\(\cos(\beta-g)=\cos\beta\), tedy \(\beta=g/2\) nebo \(\beta=\pi+g/2\),
obojí mimo \((a,\pi)\). Hodnota \(2\pi\) tedy není nikdy dosažena.
Správná strana se určí bez numeriky v limitě \(v\to1\). Pro \(0<d<\pi\)

\[
\varphi_1(d)=\pi+d+2\arctan\left(\frac14\tan\frac d2\right).
\]

Zde \(d=a-g<\pi/2\), neboť \(a<3\pi/4\), takže
\(\varphi_1(d)<\pi+2d<2\pi\). Spojitost dává tvrzení pro všechna
\(-1<v<1\).

Jedinou možnou změnou počtu kořenů na dolní hraně je \(F_L=-2\pi\). Pro
\(q=-1\)

\[
D_o=(x+1)\left[x^2-\left(\frac85+\frac25v\right)x+1\right].
\tag{BC16}
\]

První faktor je mimo rozsah \(\alpha-g\). Druhý dává
\(\cos(\alpha-g)=\cos\alpha\), tedy přesně

\[
\alpha=\frac g2,\qquad
v=v_a:=5\cos\frac g2-4=\frac{5\sqrt5-11}{4},
\qquad 0<v_a<1.
\tag{BC17}
\]

Znaménko se opět určí v koncích. Pro \(v\to1\) je \(\alpha\to0\); bod
\(\cos g-1-4i\sin g\) má argument menší než \(-\pi/2\), a proto
\(F_L<-2\pi\). Pro \(v\to-1\) položme \(d=g-b\). Platí \(0<d<b\),
protože \(b<g<2b\). Dále

\[
\varphi_{-1}(-d)
=-d-2\arctan\left(4\tan\frac d2\right)>-\pi.
\]

Poslední nerovnost je ekvivalentní \(\tan(d/2)<1/2=\tan(b/2)\). Proto
\(F_L>-2\pi\) u levého konce. Jelikož jiné nuly \(F_L+2\pi\) nejsou,

\[
\begin{cases}
F_L>-2\pi,& -1<v<v_a,\\
F_L=-2\pi,& v=v_a,\\
F_L<-2\pi,& v_a<v<1.
\end{cases}
\tag{BC18}
\]

Z (BC12)--(BC18) plyne úplný výsledek:

\begin{itemize}
\tightlist
\item
  \textbf{Hlavní kladná větev} \(\omega_c(v)\) existuje pro všechna
  \(-1<v<1\) a je jediným řešením \(F=0\).
\item
  \textbf{Dodatečná nízká větev} \(\omega_\ell(v)\) existuje právě pro
  \(v_a<v<1\) a je jediným řešením \(F=-2\pi\).
\item
  Žádná jiná celá hodnota v (BC11) neleží mezi koncovými fázemi.
\end{itemize}

\hypertarget{globuxe1lnuxed-fuxe1zovuxe9-meze-a-zuxe1vislost-na-v}{%
\subsection{Globální fázové meze a závislost na
v}\label{globuxe1lnuxed-fuxe1zovuxe9-meze-a-zuxe1vislost-na-v}}

Protože \(g\in I_+(v)\), je \(F(v,g)=-\eta(v,g)\in(-\pi,0)\). Z
monotónnosti tedy

\[
\alpha(v)<\omega_\ell(v)<g<\omega_c(v)<\beta(v).
\tag{BC19}
\]

Pro hlavní větev lze dát uniformně mnohem menší obor:

\[
\boxed{g<\omega_c(v)<g+b<a.}
\tag{BC20}
\]

Pro důkaz horní meze použijme \(\varphi_{-1}(b)=b+2\arctan2=b+a=\pi\). Z
(BC9) je pro \(-1<v<1\) hodnota \(\varphi_v(b)>\pi\). V bodě
\(\omega=g+b\) proto \(F>0\), zatímco \(F(v,g)<0\). Tento bod leží
uvnitř každé kladné mezery: nerovnost \(g+b<a\) je ekvivalentní
\(b<3\pi/10\), což plyne z

\[
\cos\frac{3\pi}{10}=\sin\frac\pi5<\frac35,
\qquad
\sin^2\frac\pi5=\frac{5-\sqrt5}{8}<\frac9{25}.
\]

Na hlavní větvi je \(\delta>0\). Přímá derivace v kladné mezeře dává

\[
F_v=
\frac{8\sin\delta}
     {(\cos\delta-v)^2+16\sin^2\delta}
+\frac1{\sqrt{16-t^2}},
\qquad
\frac{\mathrm d\omega_c}{\mathrm dv}=-\frac{F_v}{F_\omega}<0.
\tag{BC21}
\]

Hlavní větev je tedy přísně klesající ve \(v\), čili přísně rostoucí v
\(P\in(0,\pi)\).

Dodatečná větev má uniformní meze

\[
\boxed{\frac g4<\omega_\ell(v)<g\qquad(v_a<v<1).}
\tag{BC22}
\]

Stačí vyloučit \(0<\omega\le g/4=\pi/10\). Pokud takový bod leží v
mezeře, pak

\[
v>5\cos\omega-4>\frac34,
\qquad
\cos(\omega-g)\le\cos\frac{3\pi}{10}<\frac35.
\]

Použita byla přesná nerovnost \(\cos(\pi/10)>19/20\), kterou ověří
\(\cos^2(\pi/10)=(5+\sqrt5)/8>361/400\). Reálná i imaginární část
elipsového bodu jsou záporné, takže \(\theta_v(\omega-g)<-\pi/2\) a

\[
\varphi_v(\omega-g)<-\pi-\frac{3\pi}{10}.
\]

Současně

\[
\cos\eta=\frac{v-5\cos\omega}{4}
< -\frac{15}{16}
< -\cos\frac{3\pi}{10},
\qquad \eta>\frac{7\pi}{10}.
\]

Tedy \(F<-2\pi\), takže zde řešení dodatečné větve není.

\hypertarget{dux16fleux17eituxe1-nelinearita-nuxedzkuxe9-vux11btve}{%
\subsection{Důležitá nelinearita nízké
větve}\label{dux16fleux17eituxe1-nelinearita-nuxedzkuxe9-vux11btve}}

Nízká větev není monotónní ve \(v\). Toto tvrzení lze prokázat bez
vykreslení nebo numeriky. Položme \(c_*=\cos(g/2)\). Při \(\omega=g/2\)
se reálný kandidátní polynom \[
\Phi(v,\omega)
=(4c-v+5\cos\omega)(c-v)^2
-16s^2(6c-v-5\cos\omega),\qquad
c=\cos(\omega-g),\quad s=\sin(\omega-g),
\] který plyne z reálné části (BC7), přesně faktorizuje:

\[
\Phi(v,g/2)
=(c_*-v)\big[(v-5c_*)^2-16\big].
\tag{BC23}
\]

V intervalu \((v_a,1)\) je jeho jediným kořenem \(v=c_*\); tento kořen
je skutečně fyzikální: \(q=-e^{ig/2}\), \(\eta=\pi-g/2\), \(F=-2\pi\). V
tomto bodě

\[
F_v(c_*,g/2)=-\frac1{4\sin(g/2)}<0.
\tag{BC24}
\]

Protože nejsou jiné vnitřní průsečíky a \(F_\omega>0\), spojitost a
(BC24) dávají

\[
\begin{cases}
\omega_\ell(v)<g/2,& v_a<v<c_*,\\
\omega_\ell(v)=g/2,& v=c_*,\\
\omega_\ell(v)>g/2,& c_*<v<1.
\end{cases}
\tag{BC25}
\]

V limitě \(v\downarrow v_a\) přitom \(\omega_\ell\to g/2\). Má tedy
uvnitř \((v_a,c_*)\) alespoň jedno minimum. Tento krok sám ještě
nedokazuje jedinost minima. Ta plyne z kladné druhé derivace v každém
kritickém bodě, doložené níže v (L5). Pro celý soupis tedy nelze
předpokládat monotónnost nízké větve.

\hypertarget{zuxe1pornuxe1-mezera-celkovuxfd-poux10det-a-prahovuxe9-limity}{%
\subsection{Záporná mezera, celkový počet a prahové
limity}\label{zuxe1pornuxe1-mezera-celkovuxfd-poux10det-a-prahovuxe9-limity}}

Původní vzorce mají přesnou symetrii

\[
f_A(-v,-\lambda)=f_A(v,\lambda).
\tag{BC26}
\]

Skutečně \((v,\lambda,q,x)\mapsto(-v,-\lambda,-q,-x)\) ponechá \(N\)
beze změny a změní \(D_o\) na \(-D_o\). Fyzikální vrstva se zachovává.
Ekvivalentně

\[
\varphi_{-v}(\delta+\pi)=\varphi_v(\delta)+3\pi,
\qquad
\eta(-v,\omega+\pi)=\eta(v,\omega)-\pi,
\]

takže \(F(-v,\omega+\pi)=F(v,\omega)+4\pi\).

Všechny záporné póly jsou proto přesně

\[
\omega_c^-(v)=\pi+\omega_c(-v),\qquad -1<v<1,
\tag{BC27}
\]

\[
\omega_\ell^-(v)=\pi+\omega_\ell(-v),\qquad -1<v<-v_a.
\tag{BC28}
\]

Jejich celočíselné fáze jsou \(F=4\pi\), respektive \(F=2\pi\). Mají
stejné posunuté uniformní meze jako (BC20),(BC22). Dodatečná kladná a
dodatečná záporná větev nemohou být přítomny při témže \(v\).

\begin{longtable}[]{@{}
  >{\raggedright\arraybackslash}p{(\columnwidth - 4\tabcolsep) * \real{0.2500}}
  >{\raggedright\arraybackslash}p{(\columnwidth - 4\tabcolsep) * \real{0.6500}}
  >{\raggedleft\arraybackslash}p{(\columnwidth - 4\tabcolsep) * \real{0.1000}}@{}}
\toprule\noalign{}
\begin{minipage}[b]{\linewidth}\raggedright
Parametr
\end{minipage} & \begin{minipage}[b]{\linewidth}\raggedright
Přísné fyzikální A-póly
\end{minipage} & \begin{minipage}[b]{\linewidth}\raggedleft
Počet
\end{minipage} \\
\midrule\noalign{}
\endhead
\bottomrule\noalign{}
\endlastfoot
\(-1<v<-v_a\) & \(\omega_c(v),\ \pi+\omega_c(-v),\ \pi+\omega_\ell(-v)\)
& 3 \\
\(-v_a\le v\le v_a\) & \(\omega_c(v),\ \pi+\omega_c(-v)\) & 2 \\
\(v_a<v<1\) & \(\omega_\ell(v),\ \omega_c(v),\ \pi+\omega_c(-v)\) & 3 \\
\end{longtable}

Při \(v=v_a\) má dodatečná větev limit \(\omega=g/2\), \(q=-1\),
\(t=-4\); při \(v=-v_a\) je symetrickým limitem \(\omega=\pi+g/2\),
\(q=1\), \(t=4\). Tyto limity nejsou dalšími atomy A-spektrální míry. Na
obou bodech je \(\partial_\omega(t^2-16)\ne0\), \(B\ne0\), a

\[
D_o=xC+\frac t4B-\frac B4R
=-\frac{B_*}{4}R+O(R^2),\qquad R^2=t^2-16.
\tag{BC29}
\]

Proto \(f_A=O(R^{-1})=O((\lambda-\lambda_*)^{-1/2})\). Atom by vyžadoval
nenulový limit \((\lambda-\lambda_*)f_A\); tento limit je nula. V
tabulce se tedy na hranici správně počítají pouze dva přísné póly.

V původní proměnné \(P\in[-\pi,\pi]\), \(p_a=\arccos v_a\), je dodatečná
kladná větev na \(|P|<p_a\) a dodatečná záporná na \(|P|>\pi-p_a\), s
kompaktními body vyňatými podle dalšího oddílu. To je klasifikace na
celém kruhu; nepoužívá malé vybrané výstupní okno \(I_K\).

\hypertarget{pux159esnuxe9-kompaktnuxed-konce-a-nulovuxe9-kontaktnuxed-vuxe1hy}{%
\subsection{Přesné kompaktní konce a nulové kontaktní
váhy}\label{pux159esnuxe9-kompaktnuxed-konce-a-nulovuxe9-kontaktnuxed-vuxe1hy}}

Při \(v=1\) a \(v=-1\) platí původní přesná zkrácení

\[
f_A(1,\lambda)=\frac{\lambda+\gamma}{\lambda-\gamma},
\qquad
f_A(-1,\lambda)=\frac{\lambda-\gamma}{\lambda+\gamma},
\qquad \gamma=e^{-ig}.
\tag{BC30}
\]

A-kontakt tedy v prvním případě vidí pouze atom fáze \(g\) s vahou 1, ve
druhém pouze atom fáze \(\pi+g\) s vahou 1. Není správné dosadit
\(v=\pm1\) do nekráceného \(D_o\) a všechny jeho nuly prohlásit za
A-póly.

Hlavní větev má limity

\[
\lim_{v\uparrow1}\omega_c(v)=g,
\qquad
\lim_{v\downarrow-1}\omega_c(v)=\omega_H\in(g,g+b).
\tag{BC31}
\]

První limit plyne z (BC21): kdyby byl větší než \(g\), limitní
\(\varphi_1(\delta)>\pi\) by nemohl být roven \(\eta<\pi\). Druhý je
jednoznačným řešením

\[
(\omega_H-g)
+2\arctan\left(4\tan\frac{\omega_H-g}{2}\right)
=\arccos\frac{-1-5\cos\omega_H}{4},
\qquad g<\omega_H<g+b.
\tag{BC32}
\]

Dodatečná větev má limity

\[
\lim_{v\downarrow v_a}\omega_\ell(v)=g/2,
\qquad
\lim_{v\uparrow1}\omega_\ell(v)=\omega_L\in(g/2,g),
\tag{BC33}
\]

kde \(\omega_L\) je jediné řešení

\[
(g-\omega_L)
+2\arctan\left(\frac14\tan\frac{g-\omega_L}{2}\right)
+\arccos\frac{1-5\cos\omega_L}{4}=\pi.
\tag{BC34}
\]

Levá strana (BC34) je přísně klesající v \(\omega\in(0,g)\), na levém
konci větší než \(\pi\) a na pravém menší než \(\pi\). Dolní mez
\(\omega_L>g/2\) plyne také z (BC25) a z nenulové hodnoty limitního
\(F+2\pi\) při \((v,\omega)=(1,g/2)\). Tyto limitní fáze jsou navzájem
oddělené od jediného viditelného atomu.

Při \(v\uparrow1\) tudíž hlavní kladná větev přechází do jediného
viditelného atomu \(g\), zatímco póly limitních fází \(\omega_L\) a
\(\pi+\omega_H\) ztrácejí A-váhu. Při \(v\downarrow-1\) platí symetrický
výrok. Nulové váhy plynou i přímo ze slabé spojitosti A-spektrální míry:
objektový operátor je spojitý v \(P\), limitní kontaktní funkce je
(BC30) a zmíněné limitní fáze jsou od jejího atomu oddělené.

Tento výrok se týká kontaktních vah. Neříká, že všechny A-neviditelné
limitní stavy z celého objektového Hilbertova prostoru zmizí.

\hypertarget{vazba-na-stacionuxe1rnuxed-uxfalohu}{%
\subsection{Vazba na stacionární
úlohu}\label{vazba-na-stacionuxe1rnuxed-uxfalohu}}

Do plného \(F_2\) patří všechny právě klasifikované větve spolu s celým
kontinuem a oběma pulzními znaménky. Stacionární spojení splňuje přesně

\[
e=\omega_j(v)+\sigma\varepsilon(k),\qquad
\sin P\,\frac{F_v}{F_\omega}
=\sigma\varepsilon'(k),\qquad v=\cos P,\quad K=P+k.
\tag{BC35}
\]

Jde o původní neumocněnou rovnost rychlostí. Následující oddíly určují,
které členy úplného soupisu ji mohou splnit, a přímo odhadují společnou
druhou derivaci.

\hypertarget{oddux11blenuxe1-kontrola-globuxe1lnuxed-lichuxe9-b-puxf3ly}{%
\subsection{Oddělená kontrola: globální liché
B-póly}\label{oddux11blenuxe1-kontrola-globuxe1lnuxed-lichuxe9-b-puxf3ly}}

Úplný lichý kontaktní vzorec je

\[
f_B(v,\omega)=
\frac{R-4i\sin(\omega-g)}{v-5\cos\omega+4\cos(\omega-g)}.
\tag{BC36}
\]

Na přísné mezeře je jeho skutečný pól právě tehdy, když

\[
v=5\cos\omega-4\cos(\omega-g),
\qquad \sin\omega\,\sin(\omega-g)<0.
\tag{BC37}
\]

Opačné znaménko by v čitateli dalo zkrácení. Pro každé \(-1\le v\le1\)
má (BC37) právě jeden kladný kořen \(\omega_B(v)\in(0,g)\) a druhý
\(\pi+\omega_B(-v)\in(\pi,\pi+g)\). V kladném intervalu je pravá strana
(BC37) přísně klesající: její derivace je
\(-5\sin\omega+4\sin(\omega-g)<0\), na \(0\) je větší než 1 a na \(g\)
menší než \(-1\).

Dokonce lze psát

\[
R_B=\sqrt{41-40\cos g},\qquad
\vartheta=\arctan\frac{4\sin g}{5-4\cos g},
\qquad
\boxed{\omega_B(v)=\arccos\frac v{R_B}-\vartheta.}
\tag{BC38}
\]

Tyto póly nevstupují do A-resolventu \(F_2\). V původním \(F_1\) okně
\(e=4\pi/5\pm1/1000\) také žádný není, protože toto okno leží mezi \(g\)
a \(\pi\), zatímco obě liché rodiny leží mimo tento interval. Jejich
uvedení kontroluje oddělení obou úplných kontaktních sektorů; nemění
původní pravděpodobnostní jev.

\section{Hlavní kladná vázaná větev nemá stacionární spojení s pulzem}

\textbf{NON-CANONICAL / L1. Analytické odvození.}

Použijeme hlavní kladný fyzikální pól z úplného fázového soupisu.
Pro \(-1<v<1\), \(v=\cos P\), je větev \(\omega_c(v)\)
jediným řešením \(F=0\). V hybnostní proměnné ji označujeme také
\(\omega_A(P)=\omega_c(\cos P)\). Položme
\[
\omega=\omega_A(P),\qquad
\delta=\omega-g,\qquad c=\cos\delta,\qquad s=\sin\delta,
\]
\[
A=c-v,\qquad L=A^2+16s^2,\qquad
S=\sqrt{16-t^2},\qquad t=v-5\cos\omega.
\]

Z fázové rovnosti \(\delta+2\theta=\eta\in(0,\pi)\) plyne
\(\delta>0\), \(\theta>0\) a
\(\theta=(\eta-\delta)/2<\pi/2\). Tedy \(A>0\) a
\[
\frac{4s}{c-v}=\tan\theta<\cot(\delta/2)=\frac{1+c}{s}.
\]
Proto
\[
\boxed{c>\frac{v+4}{5}>v,\qquad
0<\delta<\alpha(v)\le b=\arccos(3/5).}
\tag{M1}
\]
Zejména \(g<\omega<g+b<e\) v původním energetickém okně pro
\(e\). Větev se může spojit pouze s pulzem \(\sigma=+1\).

Pro \(P\in(0,\pi)\) definujme
\[
w_1=\frac{8(1-vc)}L,\qquad w_2=\frac{5\sin\omega}S,
\qquad r_1=\frac{\sin P\,s}{1-vc},\qquad
r_2=\frac{\sin P}{5\sin\omega}.
\]
Implicitní derivace fázové rovnice dává přesně
\[
\omega_A'(P)=\frac{w_1r_1+w_2r_2}{1+w_1+w_2}.
\tag{M2}
\]
Pro pevné \(v\) má výraz
\((1-v^2)(1-c^2)/(1-vc)^2\) derivaci podle \(c\) se znaménkem
\(v-c\). Z~(M1) proto
\[
r_1^2<\frac{(1+v)(v+9)}{(v+5)^2}
=1-\frac{16}{(v+5)^2}<\frac59.
\]
Také \(\omega\in(g,g+b)\subset(\pi/4,3\pi/4)\), protože
\(b<\pi/3\) a \(g+b<11\pi/15<3\pi/4\). Tedy
\(r_2<\sqrt2/5<3/10\). Ze skutečných kladných vah v~(M2)
vychází
\[
\boxed{|\omega_A'(P)|<\frac{\sqrt5}{3}<\frac34.}
\tag{M3}
\]

Při energetickém spojení má pulz
\(\varepsilon=e-\omega<g+1/1000\). Kosinus je Lipschitzův
s konstantou jedna a \(\cos g=(\sqrt5-1)/4>307/1000\), tedy
\(\cos\varepsilon>3/10\). Skutečná pulzní rychlost proto
splňuje
\[
|\varepsilon'(k)|^2
=\frac{3+5\cos\varepsilon}{5[1+\cos\varepsilon]}
>\frac9{13}>\frac{16}{25}.
\tag{M4}
\]
Rovnost rychlostí není možná. Symetrie v~\(P\) dává totéž na
celém kruhu. Kompaktní viditelný bod \(P=0\) má nulovou
objektovou rychlost a toto vyloučení neporušuje; na opačném
kompaktním konci kontaktní váha zaniká.

\section{Nízké vázané větve: přesný obor a vyloučení stacionarity}

\textbf{NON-CANONICAL / L1. Samostatné analytické odvození,
29. října 2026.}

Tento oddíl navazuje na úplný fázový soupis \(A\)-viditelných
sudých objektových pólů. Původní rozptylový jev ani jeho
normalizace se nemění. Odvození neobsahuje numerický nebo
strojový vědecký běh.

\subsection{Vstup z fázového soupisu}

Pišme \(g=2\pi/5\), \(v=\cos P\), \(c_g=\cos g\) a
\(v_a=5\cos(g/2)-4>0\). V kladné mezeře
\[
\alpha(v)=\arccos((v+4)/5)<\omega<\arccos((v-4)/5)
\]
je vedle hlavního pólu ještě právě jedna nízká větev
\(\omega_\ell(v)\), a to přesně pro \(v_a<v<1\). Platí
\(\alpha(v)<\omega_\ell(v)<g\). Její limit při
\(v\downarrow v_a\) je \(g/2\). Příslušná záporná větev je
\(\pi+\omega_\ell(-v)\). Při derivování podle hybnosti \(P\)
se níže vždy rozumí složená funkce \(\omega_\ell(\cos P)\).

Na nízké větvi označme
\[
a=g-\omega\in(0,g),\qquad c=\cos a,\qquad s=\sin a>0,
\qquad A=c-v,\qquad L=A^2+16s^2,
\]
\[
t=v-5\cos\omega,\qquad S=\sqrt{16-t^2}>0,
\qquad \eta=\arccos(t/4),\qquad \chi=(\eta+a)/2.
\]
Se spojitým argumentem
\(\Theta=\operatorname{Arg}(c-v+4is)\in(0,\pi)\) je
fyzikální rovnice
\[
F(v,\omega)=-a-2\Theta-\eta=-2\pi,
\qquad H=A\sin\chi+4s\cos\chi=0.
\tag{L1}
\]
Přitom \(H=-\sqrt L\sin(F/2)\), \(-4\pi<F<0\) a
\[
F_\omega=1+\frac{8(1-vc)}L+\frac{5\sin\omega}S>0,
\qquad F_v=-\frac{8s}L+\frac1S.
\tag{L2}
\]
Proto \(H<0\) znamená \(F<-2\pi\) a \(H>0\) znamená
\(F>-2\pi\). Fyzikální nízký kořen leží napravo, respektive
nalevo od takového testu.

\subsection{Každý vnitřní kritický bod je přísné minimum}

Při \(F_v=0\) je \(L=8sS\). Přímá derivace při pevném
\(\omega\) dává
\[
F_{vv}=-\frac{16sA}{L^2}+\frac{t}{S^3}
=\frac{t-AS/(4s)}{S^3}.
\tag{L3}
\]
Z \(\omega<g\), \(v\le1\) a \(5\cos g>1\) plyne
\(t<0\). Je-li \(A\ge0\), je čitatel záporný. Je-li
\(A<0\), použijeme \(-A\le1-c\) a \(S\le4\):
\[
\frac{AS}{4s}\ge-\frac{1-c}{s}=-\tan(a/2).
\]
Konkávnost kosinu a konvexita tangenty na příslušných intervalech
dávají
\[
\cos(g-a)\ge c_g+(1-c_g)a/g,
\qquad \tan(a/2)\le(a/g)\tan(g/2).
\]
Proto
\[
5\cos\omega-1-\tan(a/2)
\ge5c_g-1+\frac ag[5(1-c_g)-\tan(g/2)]>0.
\tag{L4}
\]
Stačí \(c_g>1/5\), \(c_g<1/3\) a \(\tan(g/2)<1\).
Tedy také v tomto případě \(t-AS/(4s)<0\). Z implicitní
derivace při \(\omega_\ell'(v)=0\) plyne
\[
\omega_\ell''(v)=-F_{vv}/F_\omega>0.
\tag{L5}
\]
Každý vnitřní kritický bod je přísné minimum. Nemůže jich být
více než jeden, protože mezi dvěma přísnými minimy by existovalo
vnitřní maximum. Horní mez celé větve proto stačí kontrolovat
na jejích koncích.

\subsection{Horní mez nízké fáze}

Na prahovém konci je, v označení spojitého limitu,
\(\omega_\ell(v_a)=g/2=\pi/5<2/3\). Na druhém konci
\(v=1\) testujme \(\omega=2/3\). Tento bod je uvnitř
kladné mezery. Položme
\[
r=\tan((g-2/3)/2),\qquad
y=\tan(\eta/2)
=\sqrt{\frac{3+5\cos(2/3)}{5-5\cos(2/3)}}.
\]
Z \(\pi<22/7\) je \((g-2/3)/2<31/105\). Současně
\[
\arctan(61/200)>\frac{61}{200}
-\frac13\left(\frac{61}{200}\right)^3
=\frac{7093019}{24000000}>\frac{31}{105}.
\]
Tedy \(r<61/200\). Dále
\[
\cos(2/3)\le1-\frac{(2/3)^2}{2}+\frac{(2/3)^4}{24}
=\frac{191}{243}<\frac{2361}{3001},
\]
což je právě podmínka \(y<51/20\). Při \(v=1\) má \(H\)
přesně znaménko \(4-r^2-5ry\); všechny vytažené faktory jsou
kladné. Nyní
\[
r^2+5ry<\frac{3721}{40000}+\frac{15555}{4000}
=\frac{159271}{40000}<4.
\tag{L6}
\]
Proto \(H(1,2/3)>0\). Kořen \(F=-2\pi\) má tedy při
\(v\to1\) limit menší než \(2/3\). Tento argument zároveň
brání jeho úniku ke kompaktnímu pólu \(\omega=g\). Kdyby
někde uvnitř větev dosáhla \(2/3\), měla by mezi dvěma nižšími
konci vnitřní maximum, což odporuje~(L5). Platí
\(\omega_\ell<2/3\) všude.

\subsection{Dolní mez nízké fáze}

Testujme \(\omega=1/2\). Potřebné meze jsou
\[
a=g-1/2>3/4,\qquad c<3/4,\qquad s>2/3,
\qquad u=\cos(1/2)>7/8.
\tag{L7}
\]
Použita byla \(\pi>25/8\). Tuto slabou racionální mez lze
doložit například Machinovou identitou:
\[
\pi=16\arctan(1/5)-4\arctan(1/239)
>\frac{1184}{375}-\frac4{239}
=\frac{281476}{89625}>\frac{25}{8}.
\]
Zbylé nerovnosti plynou z prvních střídavých členů sinu a kosinu.

Je-li \(\alpha(v)\ge1/2\), je závěr okamžitý. Jinak test
leží v mezeře a \(v>5u-4>3/8\).

Pro \(v\le3/5\) máme
\(\cos\eta<(3/5-35/8)/4=-151/160\). Protože
\(\cos(3/8)<15/16<151/160\), je \(\eta>\pi-3/8\).
Tedy \(\chi>\pi/2+3/16\) a
\(-\cos\chi>\sin(3/16)>9/50\). Odtud
\[
H<\frac38-\frac83\frac9{50}<0.
\]
Pro \(v\ge3/5\) máme \(A<3/20\) a
\(\cos\eta<(1-35/8)/4=-27/32\). Protože
\(\cos(3/5)<27/32\), je \(\eta>\pi-3/5\). Tedy
\(\chi>\pi/2+3/40\) a
\(-\cos\chi>\sin(3/40)>7/100\). Odtud
\[
H<\frac3{20}-\frac83\frac7{100}<0.
\]
V obou případech \(\chi<\pi\), takže použité znaménko kosinu
a horní mez prvního členu jsou správné. Podle~(L1),(L2) leží
nízký kořen nad testem. Celkově
\[
\boxed{\frac12<\omega_\ell(v)<\frac23\qquad(v_a<v<1).}
\tag{L8}
\]

\subsection{Rychlost celé nízké větve je menší než jedna polovina}

Pro \(P\in(0,\pi)\) pišme
\[
w_1=\frac{8(1-vc)}L,\qquad w_2=\frac{5\sin\omega}S,
\qquad r_1=\frac{\sin P\,s}{1-vc},\qquad
r_2=\frac{\sin P}{5\sin\omega}.
\]
Přesně
\[
\frac{\mathrm d}{\mathrm dP}\omega_\ell(\cos P)
=\frac{-w_1r_1+w_2r_2}{1+w_1+w_2}.
\tag{L9}
\]
Máme \(0\le r_1\le1\), protože
\((1-vc)^2-(1-v^2)s^2=(v-c)^2\). Z~(L8) je
\(a>g-2/3>7/12\) a tedy \(s>11/20\). Také
\(v>v_a>0\) a \(c>0\). Přesná identita
\[
(1-vc)(c^2+16s^2)-L
=v[c(-14+15c^2)-v]\le0
\]
dává, protože již doložené \(s>11/20\) znamená
\(c^2<279/400<14/15\),
\[
w_1\le\frac8{1+15s^2}<\frac{640}{443},\qquad
w_2>\frac54\sin(1/2)>\frac{115}{192},\qquad
r_2<\frac{48}{115}<\frac12.
\tag{L10}
\]
Zde \(640/443<1+115/192\). Kladná část~(L9) je menší než
\(r_2<1/2\). Záporná část má modul menší než
\(w_1/(1+w_1+w_2)<1/2\). Po symetrii v~\(P\) platí
\[
\boxed{\left|\frac{\mathrm d}{\mathrm dP}
\omega_\ell(\cos P)\right|<\frac12.}
\tag{L11}
\]

\subsection{Obě nízké větve jsou v původním okně nestacionární}

Původní \(e\) leží v intervalu
\([4\pi/5-1/1000,4\pi/5+1/1000]\). Pro kladnou nízkou
větev je možná jen \(\sigma=+1\) a pulzní fáze je
\(\varepsilon=e-\omega_\ell\). Pak
\[
\pi-\varepsilon>\pi/5-1/1000+1/2>109/100.
\]
Přesná kosinová mez
\[
\cos(109/100)\le
\frac{1115438161}{2400000000}<\frac7{15}
\]
dává \(\cos\varepsilon>-7/15\). Ze skutečné pulzní rychlosti
\[
|\varepsilon'(k)|^2
=\frac{3+5\cos\varepsilon}{5[1+\cos\varepsilon]}
\]
proto plyne \(|\varepsilon'|>1/2\). Rovnost rychlostí
s~(L11) není možná.

Pro \(\pi\)-posunutou nízkou větev je možná jen
\(\sigma=-1\). Její pulzní fáze splňuje
\[
0<\varepsilon=\pi+\omega_\ell-e
<\pi/5+2/3+1/1000<\pi/2.
\]
Tedy \(|\varepsilon'|^2>3/5>1/4\), zatímco objektová
rychlost má po symetrii stále modul menší než \(1/2\). Ani
tato větev nemá žádný stacionární bod v původním energetickém
okně.

Prahové konce se posuzují společným pólovým a kontinuálním
příspěvkem; kompaktní konce s nulovou \(A\)-váhou nejsou
započítávány jako nové viditelné póly.

\section{Zbývající fyzikální stacionární větve vázaných pólů}
\label{sec:bound-stationary}

\textbf{NON-CANONICAL / L1. Analytický pracovní důkaz, 29. října 2026.}
Nebyl použit numerický průzkum kořenů, vyhodnocení CAS, vědecký běh ani
veřejná změna repozitáře. Všechny následující odhady jsou symbolické
identity a racionální nerovnosti. Zachovává se celý původní kruh \(K\),
původní energetické okno, normalizace zdroje i pravděpodobnostní jev.

Tento oddíl řeší hlavní vázanou větev posunutou o \(\pi\), se záporným
pulzním znaménkem. Doprovodný fázový soupis a rychlostní vyloučení ji
určují jako jedinou zbývající možnou stacionární vázanou rodinu:
hlavní větev v horní energetické půlperiodě a obě nízké větve již byly
vyloučeny rychlostními mezemi. Zde dokážeme existenci právě dvou
fyzikálních stacionárních křivek a jednotnou nenulovou společnou druhou
derivaci.

Východiskem jsou fyzikální jmenovatel (BC6), fázová rovnice
(BC11) a původní úplné kontaktní vzorce [S1,S2]. Společná regularita
pólu a hrany v jiných integračních bodech je samostatná úloha; tento
oddíl netvrdí globální analytičnost \(F_2\).

\subsection{Fyzikální půlperiodová symetrie a příslušná větev}

Položme
\[
g=2\pi/5,\qquad C=\cos g,\qquad S_g=\sin g,
\qquad a_p=\arccos(-3/5),
\]
\[
e\in I_e=[4\pi/5-10^{-3},4\pi/5+10^{-3}].
\]
Úplný sudý kontaktní výraz je invariantní při záměně
\[
(v,\omega,x,q)\longmapsto(-v,\omega+\pi,-x,-q).
\tag{BS1}
\]
Jeho čitatel \(N\) se totiž nezmění, jmenovatel \(D_o\) změní znaménko
a podíl \(2xN/D_o\) zůstane stejný. Správně se přenáší také fyzikální
mezerový kořen: \(t\) i předepsaná hraniční hodnota \(R\) změní
znaménko, a proto je změní i \(q\). Kontaktní spektrální váha \(A\)
se při (BS1) zachovává.

Hlavní vázanou fázi v horní půlperiodě pišme jako
\[
\omega_A(p)=g+\delta(\cos p),\qquad 0\le p\le\pi,
\qquad \delta>0\quad\text{pro }p\ne0.
\]
V \(p=0\) je její fází původní kompaktní vlastní hodnota \(g\).
Partner posunutý o \(\pi\) má pro \(P=\pi-p\) fázi
\(\pi+\omega_A(p)\). Stacionární rezonance v původním okně \(e\)
může použít pouze záporné pulzní znaménko:
\[
e=\pi+g+\delta-\varepsilon(k),\qquad
\varepsilon(k)=d(e)+\delta,
\qquad d(e)=\pi+g-e.
\tag{BS2}
\]
Přitom
\[
3\pi/5-10^{-3}\le d(e)\le3\pi/5+10^{-3},
\qquad 0<d(e)<a_p.
\tag{BS3}
\]
Poslední nerovnost plyne již z \(a_p>2\pi/3\). Jediný přípustný
kladný interval parametru je tedy
\(0\le\delta\le\delta_{\max}(e):=a_p-d(e)\).

\subsection{Výslovná jednoparametrická reprezentace na fyzikální vrstvě}

Označme
\[
c=\cos\delta,\qquad s=\sin\delta,
\qquad \mathfrak a=5(1-C),\qquad \mathfrak b=3+5C,
\]
\[
\mathfrak h=5S_g,
\qquad \mathfrak a+\mathfrak b=8.
\]
Použijme parametr kontaktního úhlu
\[
z=\frac{c-v}{4s}>0,\qquad
\theta=\arctan(1/z),\qquad \eta=\delta+2\theta.
\tag{BS4}
\]
Reálná kontaktní rovnice \(H=0\) je pak splněna identicky, protože
\(c-v=4sz=4s\cot\theta\). Zbývající volný mezerový vztah
\(v-5\cos(g+\delta)=4\cos\eta\) je přesně
\[
\boxed{\mathscr G(z,\delta)=cU(z)-sV(z)=0,
\qquad v=c-4sz,}
\tag{BS5}
\]
kde
\[
U(z)=\mathfrak b z^2-\mathfrak a,\qquad
V(z)=\mathfrak h(1+z^2)+4z-4z^3.
\tag{BS6}
\]
Například rozdíl obou stran volného mezerového vztahu, vynásobený
\(1+z^2\), je právě \(-\mathscr G\). Jde o původní jmenovatel
a fyzikální kořen, nikoli o výběr z umocněné rovnice.

Pro \(0<\delta\le1/3\) má rovnice právě jeden kladný kořen \(z\).
Funkce
\[
\frac{\mathscr G}{z^2}
=4sz+(\mathfrak b c-\mathfrak h s)
-\frac{4s}{z}-\frac{\mathfrak a c+\mathfrak h s}{z^2}
\]
je totiž striktně rostoucí od minus nekonečna do plus nekonečna.
Její kladný kořen je jednoduchý a hladce závisí na \(\delta\).
Položme
\[
z_0=\sqrt{\mathfrak a/\mathfrak b}.
\]
Elementární přesné meze
\[
\frac{309}{1000}<C<\frac{31}{100},\qquad
\frac{19}{20}<S_g<\frac{119}{125}
\tag{BS7}
\]
dávají \(6/7<z_0<1\). K ověření (BS7) stačí umocnit racionální
meze \(559/250<\sqrt5<56/25\) a použít \(S_g^2=1-C^2\).
V \(z=z_0\) je \(U=0\), \(V>0\), tedy \(\mathscr G<0\)
a kladný kořen splňuje \(z>z_0\).

Platí také jednotná horní mez \(z<5/4\). Přímo
\[
16\mathscr G(5/4,\delta)
=c(205C-5)+s(45-41\mathfrak h).
\]
Protože \(c\ge17/18\), \(s\le1/3\), dostáváme z (BS7)
\[
16\mathscr G(5/4,\delta)
>\frac{17}{18}\frac{11669}{200}
-\frac13\frac{3754}{25}
=\frac{18181}{3600}>0.
\tag{BS8}
\]
V celém tomto parametrickém intervalu proto
\[
\boxed{6/7<z<5/4,\qquad -13/18<v<1.}
\tag{BS9}
\]
Dolní mez \(v\) plyne z
\(v=c-4sz>17/18-5/3=-13/18\). Navíc
\(\tan(\delta/2)\le12/71<6/7<z\), takže (BS4) dává
\(0<\eta<\pi\). Rovnice (BS5) tedy leží na fyzikální kladné
mezerové vrstvě, při \(-1<v<1\), a jde o jednoduchý viditelný pól.
Při \(\delta=0\) má spojitý konec \(v=1\), \(z=z_0\), s fází
\(g\). Celý interval \([0,1/3]\) je tedy na původní hlavní větvi
skutečně dostupný.

Z \(\tan\delta=U/V\) lze též získat úplně výslovnou algebraickou
parametrizaci:
\[
c=\frac{V}{\sqrt{U^2+V^2}},\qquad
s=\frac{U}{\sqrt{U^2+V^2}},\qquad
v=\frac{V-4zU}{\sqrt{U^2+V^2}}.
\tag{BS10}
\]
Na tomto úseku jsou \(U,V>0\). Kladná odmocnina a fyzikální
úhlový interval v (BS4) jednoznačně určují větev.

\subsection{Energetická podmínka vynutí parametr menší než třetina}

Nechť \(\Delta=a_p-3\pi/5\), takže \(0<\Delta<\pi/2\).
Přesné odečtení úhlů dává
\[
\tan\Delta=\frac{3S_g-4C}{3C+4S_g}.
\]
Z (BS7) je \(C/S_g>309/952\). Zobrazený racionální výraz s tímto
poměrem klesá, a tedy
\[
\tan\Delta<\frac{324}{947}<\frac{343}{1000}.
\]
Pro \(x=997/3000=1/3-10^{-3}\) platí
\(\tan x\ge x+x^3/3>343/1000\). První nerovnost plyne integrací
\((\tan x)'=1+\tan^2x\ge1+x^2\); druhá je racionální porovnání,
pro které již stačí \(x>33/100\). Proto
\[
\boxed{0<\delta_{\max}(e)\le\Delta+10^{-3}<1/3.}
\tag{BS11}
\]
Tím je každý možný stacionární bod této větve lokalizován ještě před
použitím derivací.

\subsection{První rychlostní mez vynutí parametr větší než čtvrtina}

Pišme
\[
S=\sqrt{16-(v-5\cos\omega)^2}=4\sin\eta>0,
\qquad L=(c-v)^2+16s^2,
\qquad \omega=g+\delta.
\]
Derivace přesné fázové rovnice
\(\delta+2\operatorname{Arg}(c-v+4is)-\eta=0\) dávají
\[
\omega_v=-\frac{N}{D_*},\qquad N=L+8sS,
\qquad D_*=5\sin\omega\,L+S[L+8(1-vc)].
\tag{BS12}
\]
Veličiny \(N,D_*\) jsou kladné. Zvláště \(v(\delta)\) striktně
klesá. Pro \(0<\delta\le1/3\) je
\(\sin\omega\ge S_g>19/20\): použili jsme
\(g\le\omega\le g+1/3<\pi-g\). Současně
\[
1-vc=s(s+4cz)>3s,
\]
protože \(c>7/8\), \(z>6/7\). Proto
\[
D_*-3N=(5\sin\omega-3)L+S[L+8(1-vc-3s)]>0.
\]
Odtud
\[
0<\omega_A'(p)=\sqrt{1-v^2}(-\omega_v)<1/3.
\tag{BS13}
\]
Ve stacionární rezonanci se tato rychlost rovná velikosti pulzní
rychlosti. Pro \(\varepsilon=d(e)+\delta\) tedy přesný pulzní vztah
dává
\[
|\varepsilon'|^2=
\frac{3+5\cos\varepsilon}{5(1+\cos\varepsilon)}<1/9,
\qquad \cos\varepsilon<-11/20.
\]
Na druhé straně
\[
\cos(3\pi/5+251/1000)
>-\frac{31}{100}-\frac{119}{125}\frac{251}{1000}
=-\frac{68619}{125000}>-\frac{11}{20}.
\]
Použili jsme \(\cos x\le1\), \(\sin x\le x\) v přesném součtovém
vzorci. Všechny úhly leží v \((0,\pi)\), kde kosinus klesá.
Protože \(d(e)\le3\pi/5+10^{-3}\), musí každý stacionární bod
splňovat
\[
\boxed{1/4<\delta<\delta_{\max}(e)<1/3.}
\tag{BS14}
\]

\subsection{Užší parametrický obor a meze inverzních derivací}

Ve stacionárním bodě je \(z>1\). Jinak (BS5), (BS6) a
\(z_0<z\le1\) dávají
\[
\tan\delta=U/V
\le(\mathfrak b-\mathfrak a)/\mathfrak h
<\frac{11/10}{19/4}=\frac{22}{95}<1/4,
\]
což odporuje (BS14). Spolu s (BS8) jsme dokázali
\[
\boxed{1/4<\delta<1/3,\qquad 1<z<5/4.}
\tag{BS15}
\]
Následující odhady platí na celém tomto parametrickém oboru,
nejen na kořeni stacionární rovnice.

Platí zde
\[
s+4cz>s+4c>4+\delta-2\delta^2-\delta^3/6>4.
\]
Nahradíme-li v porovnání za (BS12) číslo 3 číslem 4, dostaneme
\[
\boxed{0<-\omega_v<1/4,\qquad v'(\delta)<-4.}
\tag{BS16}
\]
Nyní omezíme druhou inverzní derivaci bez předpokladu konvexity
objektové větve. Derivace \(V\) jsou
\[
V'=2\mathfrak h z+4-12z^2,
\qquad V''=2\mathfrak h-24z<-14.
\]
Pro \(1<z<5/4\) odtud plyne
\[
-7/2<V'<38/25,\qquad V<10.
\tag{BS17}
\]
Pro druhé tvrzení použijeme Taylorovu nerovnost s \(V''<-14\):
\(V\le2\mathfrak h+(2\mathfrak h-8)(z-1)-7(z-1)^2\).
Její horní mez je menší než
\(238/25+(38/25)^2/28<10\). Dolní mez \(V'\) plyne již
z dosazení \(z=5/4\) do této klesající funkce při
\(\mathfrak h>19/4\).

Na kořeni (BS5) je
\[
\mathscr G_z=cU'-sV'
>\frac{17}{18}\frac{909}{100}-\frac13\frac{38}{25}>8,
\]
\[
z'=\frac{V}{c\mathscr G_z}<\frac{45}{34}<4/3.
\tag{BS18}
\]
Derivace je kladná. S \(\mathfrak b<23/5\) navíc
\[
|sU'+cV'|<15/2,
\qquad
0<\mathscr G_{zz}=2\mathfrak b c+s(24z-2\mathfrak h)<17.
\]
Kladnost posledního výrazu plyne také přímo z \(2\mathfrak h<24z\).
Protože \(\mathscr G_{\delta\delta}=-\mathscr G=0\) na jeho nule,
implicitní derivování dává přesnou formuli
\[
z''=\frac{2(sU'+cV')z'-\mathscr G_{zz}(z')^2}{\mathscr G_z}.
\]
Předchozí konstanty tedy dokazují
\[
\boxed{|z''|<\frac{2(15/2)(4/3)+17(4/3)^2}{8}
=113/18<7.}
\tag{BS19}
\]
Dvě derivace \(v=c-4sz\) nyní dávají
\[
v''=-c+4sz-8cz'-4sz''>-1-8(4/3)-4(1/3)7=-21.
\]
Spolu s \(v'<-4\) je inverzní derivace zdola omezena:
\[
\omega_{vv}=-\frac{v''}{(v')^3}>-21/64.
\]
Je-li \(v''<0\), dělí se jeho dolní mez kladným
\(|v'|^3>64\); je-li \(v''\ge0\), je tvrzení okamžité.
Objektová křivost, která nemusí být kladná, tudíž splňuje
\[
\boxed{\omega_A''(p)
=-v\omega_v+(1-v^2)\omega_{vv}>-37/64.}
\tag{BS20}
\]

\subsection{Jednotná nedegenerovanost každé stacionární rezonance}

Ve stacionárním bodě označme
\(w=\omega_A'(p)=|\varepsilon'(k)|\). Z (BS16) je \(w<1/4\).
Přesnou druhou pulzní derivaci lze vyjádřit jen jeho rychlostí:
\[
-\varepsilon''(k)
=\frac{1-w^2}{2}\sqrt{4-5w^2}
>\frac{15\sqrt{59}}{128}>7/8.
\tag{BS21}
\]
Druhá derivace při pevném původním \(K\) je tedy
\[
\boxed{\begin{aligned}
\partial_P^2[\omega_{\rm lower}(P)-\varepsilon(K-P)]
&=\omega_A''(p)-\varepsilon''(k)\\
&>19/64>1/4.
\end{aligned}}
\tag{BS22}
\]
Jde o jednotně nedegenerované minimum. Tento závěr neplyne jen
z jednoduchosti objektového pólu: obě druhé derivace byly
výslovně omezeny.

\subsection{Existence a jednoznačnost včetně neumocněných znamének}

Pro pevné \(e\) vezměme \(p(\delta)=\arccos v(\delta)\) na
\([0,\delta_{\max}(e)]\). Funkce striktně roste, začíná v nule
a podle (BS9) končí před \(\pi\). Nechť \(W(\varepsilon)\)
označuje kladnou velikost pulzní rychlosti pro
\(0<\varepsilon<a_p\), se spojitou hodnotou \(W(a_p)=0\).
Definujme rozdíl
\[
\mathcal F_e(p)=\omega_A'(p)
-W(\pi+\omega_A(p)-e).
\tag{BS23}
\]
V \(p=0\) je záporný: objektová rychlost mizí, ale
\(0<d(e)<a_p\). Na druhém konci pulzní rychlost mizí a objektová
rychlost je přísně kladná, takže rozdíl je kladný. Existuje alespoň
jedna vnitřní nula.

Uvnitř intervalu je
\(W'(\varepsilon)=\varepsilon''(k)/W(\varepsilon)\).
V každé nule (BS23) se obě kladné rychlosti rovnají, a proto
\[
\mathcal F_e'(p)=\omega_A''(p)-\varepsilon''(k)>1/4.
\]
Každá nula je tedy přísným vzestupným průchodem. Dvě různé nuly
by mezi sebou vynutily sestupný průchod nebo nulu s nekladnou
derivací. Existuje proto právě jedna nula \(p_*(e)\), s jediným
\(\delta_*(e)\) v (BS14). Věta o implicitní funkci navíc dává
hladkou závislost kořene na \(e\).

Kořen lze definovat také bez výběru znamének hybností jako jediné
řešení
\[
\frac{1-v(\delta)^2}{v'(\delta)^2}
=\frac{3+5\cos(d(e)+\delta)}{5[1+\cos(d(e)+\delta)]},
\qquad 1/4<\delta<\delta_{\max}(e),
\tag{BS24}
\]
kde \(v(\delta)\) je jednoznačná fyzikální parametrizace (BS5).
Rovnost kladných velikostí rychlostí je zde automatická; které
hybnostní dvojice skutečně nastávají, však nadále určují původní
neumocněná znaménka.

Položme
\[
\kappa_*(e)=\arccos
\frac{1-5\cos(d(e)+\delta_*(e))}{4}>0.
\]
Pro kladného reprezentanta původního \(P\) je \(P=\pi-p_*\)
a objektová derivace je záporná. Záporné pulzní znaménko proto
vyžaduje \(k=-\kappa_*\). Druhý reprezentant vznikne reflexí.
Skutečná stacionární řešení jsou tedy právě
\[
\boxed{(P,k,K)=(\pi-p_*,-\kappa_*,\pi-p_*-\kappa_*),}
\]
\[
\boxed{\begin{gathered}
(P,k,K)=(-\pi+p_*,\kappa_*,-\pi+p_*+\kappa_*),\\
\sigma=-1.
\end{gathered}}
\tag{BS25}
\]
Z umocněné rovnice nevzniká žádná další znaménková kopie. Ani jeden
kompaktní hybnostní konec není stacionární rezonancí: při
\(\delta=0\) zůstává nenulová pulzní rychlost odlišná od nulové
objektové rychlosti; druhý konec leží mimo přípustný interval
\(\delta\).

\subsection{Hrubá geometrická poloha a přísně kladné váhy}

V kořeni dávají \(\sin\delta>95/384\), \(z>1\) a \(c<1\)
meze
\[
-13/18<v<1/96.
\tag{BS26}
\]
Proto \(\pi/3<p_*<5\pi/6\). Rychlostní mez \(w<1/4\) znamená
\(\cos\kappa_*>29/30\), a tedy \(\kappa_*<1/3\), jak plyne
z elementární horní Taylorovy meze pro \(\cos(1/3)\). Odtud
\[
\boxed{K_\pm(e)=\pm K_*(e),\qquad
1/6<K_*(e)=\pi-p_*-\kappa_*<2\pi/3.}
\tag{BS27}
\]
Obě křivky jsou tedy různé a jednotně oddělené od \(K=0\).
Oddělení od běžných kontinuálních stacionárních křivek je
doloženo samostatně v následujícím oddílu.

Pro samotné tyto póly dostáváme také přímý odstup od hran mezery.
Z \(1<z<5/4\) leží \(\theta=\arctan(1/z)\) mezi \(\pi/6\)
a \(\pi/4\). Rovnice (BS15) proto dává
\[
\pi/3<\eta<2\pi/3,
\qquad \boxed{|t|=4|\cos\eta|<2.}
\tag{BS28}
\]
Žádný stacionární vázaný pól tedy sám neleží ve styku pólu
s hranou. Tato vlastnost se při fyzikální půlperiodové symetrii
zachovává.

Stejné derivace navíc dávají
\(|v'|<1/3+5+16/9=64/9\). Protože z (BS26) plyne
\(\sin p>2/3\), je
\[
\boxed{3/32<w<1/4.}
\tag{BS29}
\]
Obálková derivace na kladné křivce \(K\) je \(\mathrm de/\mathrm dK=-w\).
Tedy
\[
\boxed{4<|K_\pm'(e)|<32/3.}
\tag{BS30}
\]

Také kontaktní reziduum je oddělené od nuly. Použijme původní
reálný čitatel a jmenovatel v horním hlavním pólu:
\[
H=(c-v)\sin\theta-4s\cos\theta,
\qquad G_c=4(c+v)\cos\theta+s\sin\theta.
\]
Je \(c+v>2/9\), \(\cos\theta>2/3\), takže \(G_c>16/27\).
Při pevném \(v\) dává (BS28)
\(\lvert\theta_\omega\rvert=(1+5\sin\omega/S)/2<4/3\).
Protože \(c-v=4sz<5/3\), přímé derivování \(H\) dává
\[
|H_\omega|<1/3+(5/3)(4/3)+4+(4/3)(4/3)=25/3.
\]
Všechny členy této derivace zde mají záporné znaménko, tedy
\(H_\omega<0\). Původní reziduová formule a (BS1) dokazují
na obou skutečných dolních stacionárních pólech
\[
\boxed{r_b=-\frac{G_c}{2H_\omega}>8/225>1/32,
\qquad h_\sigma=(1-\sigma\varepsilon')/2>3/8.}
\tag{BS31}
\]
Jde o skutečné viditelné stacionární póly, nikoli o algebraické
kandidáty s nulovou vahou. Jejich lokální foldové příspěvky proto
musí zůstat v úplném \(F_2\).

\subsection{Přesný rozsah klasifikace}

Spolu s doprovodným fázovým soupisem a rychlostním vyloučením
ostatních větví ponechává tento důkaz v původním energetickém okně
právě dvě fyzikální stacionární vázané rodiny (BS25). Jejich
společná křivost je jednotně větší než \(1/4\), samotné póly mají
v proměnné \(t\) odstup větší než dva od obou objektových hran
a kontaktní i pulzní váhy jsou přísně kladné s mezemi (BS31).

Tím je uzavřen dosud chybějící počet stacionárních vázaných větví
a jejich nedegenerovanost. Tento výsledek sám nepřiřazuje celému
přenosovému sloupci globální Fourierovu konstantu. Zbývající
kontinuální příspěvky, setkání pólu s hranou a výměny pulzních
bank se stále musí spojit podle vnějšího předpisu. Původní ostré
energetické konce nadále vyžadují zaplacenou normovou aproximaci
nebo samostatný oscilační odhad. Žádná dřívější konstanta
konečného času se tímto oddílem nemění.

\section{Číselné oddělení vázaných a kontinuálních stacionárních křivek}
Na celém původním energetickém okně dokážeme kruhový rozestup
větší než \(1/300\). Z předchozího návrhu [S4, rovnice PT7--PT17]
přejímáme úplnou klasifikaci dvou běžných kontinuálních hran.
Jejich fáze je jediným řešením
\[
e=\omega+\arccos\frac{\cos\omega-3}{4},\qquad 0<\omega<b,\qquad v_c(e)=5\cos\omega(e)-4.
\]
Příslušná společná druhá derivace je menší než \(-2\sqrt3/5\),
a úplný objektový jmenovatel má modul větší než \(12/25\).
Následující důkaz rozestupu používá pouze jejich uvedené polohy,
přesný střed a již doloženou variaci.

\hypertarget{pux159esnuxe9-pux159iux159azenuxed-obou-kladnuxfdch-reprezentantux16f}{%
\subsection{Přesné přiřazení obou kladných
reprezentantů}\label{pux159esnuxe9-pux159iux159azenuxed-obou-kladnuxfdch-reprezentantux16f}}

Podle předchozího oddílu, (BS25), je u vázaného pólu \(p=\arccos v\),
ale \textbf{původní} kladná objektová hybnost je
\(P=\pi-p=\arccos(-v)\). Neumocněná stacionární rovnice při
\(\sigma=-1\) určuje \(k=-\kappa_b\), a tedy

\[
K_b=P-\kappa_b=\arccos(-v)-\kappa_b.
\tag{CS1}
\]

Platí doložené meze \(v>-13/18\) a \(w=|\varepsilon'(\kappa_b)|<1/4\),
kde \(0<\kappa_b<\pi\). Druhá křivka je přesně \(-K_b\).

Podle předchozího časového návrhu {[}S4, oddíl 14{]},
(PT8),(PT9),(PT12), má kladná kontinuální křivka

\[
K_c=p_c-\kappa_c>0,\qquad
p_c=\arccos v_c(e),\qquad
\cos\kappa_c=\frac{15-v_c(e)}{16}.
\tag{CS2}
\]

Zde je \(\sigma=+1\), ale neumocněná rovnice rovněž určuje
\(k=-\kappa_c\). Její druhá křivka je \(-K_c\). Veličiny \(v\) a \(v_c\)
tedy nejsou kosiny téže původní hybnosti; znaménko v (CS1) je
podstatné.

\hypertarget{dolnuxed-mez-vuxe1zanuxe9-kux159ivky}{%
\subsection{Dolní mez vázané
křivky}\label{dolnuxed-mez-vuxe1zanuxe9-kux159ivky}}

Střídavý Taylorův odhad pro kosinus dává přesně

\[
\cos\frac34
\ge1-\frac12\left(\frac34\right)^2
+\frac1{24}\left(\frac34\right)^4
-\frac1{720}\left(\frac34\right)^6
=\frac{239759}{327680}>\frac{13}{18}.
\tag{CS3}
\]

Protože \(-v<13/18\), monotónnost arkuskosinu implikuje
\(P=\arccos(-v)>3/4\).

Přesná pulzní rychlost splňuje

\[
w^2=\frac{4\sin^2(\kappa_b/2)}
{1+4\sin^2(\kappa_b/2)}<\frac1{16}.
\]

Odtud

\[
\sin\frac{\kappa_b}{2}<\frac1{2\sqrt{15}}
<\frac5{38}
<\frac{1346}{10125}
=\frac2{15}-\frac16\left(\frac2{15}\right)^3
<\sin\frac2{15}.
\tag{CS4}
\]

První racionální porovnání se ověří umocněním \(1444<1500\); druhé je
\(50625<51148\). Sinus je na \([0,\pi/2]\) přísně rostoucí. Proto
\(\kappa_b<4/15\), a tak

\[
\boxed{K_b>\frac34-\frac4{15}=\frac{29}{60}.}
\tag{CS5}
\]

Rovněž \(K_b<P<\pi\), což později ošetří kruhový šev.

\hypertarget{hornuxed-mez-kontinuuxe1lnuxed-kux159ivky-pouze-z-doloux17eenuxe9ho-pt-vuxfdsledku}{%
\subsection{Horní mez kontinuální křivky pouze z doloženého PT
výsledku}\label{hornuxed-mez-kontinuuxe1lnuxed-kux159ivky-pouze-z-doloux17eenuxe9ho-pt-vuxfdsledku}}

Použijme původní přesný střed a doloženou variaci z (PT14) a odstavce
následujícího po (PT17):

\[
v_c(e_0)=\frac{110\varphi-159}{31},\qquad
|v_c(e)-v_c(e_0)|<\frac1{200},\qquad
e_0=4\pi/5,\quad \varphi=\frac{1+\sqrt5}{2}.
\tag{CS6}
\]

Protože \((559/250)^2=312481/62500<5\), platí
\(\varphi>809/500=1618/1000\). Tudíž na celém původním okně

\[
v_c(e)>
\frac{110(809/500)-159}{31}-\frac1{200}
=\frac{753}{1240}>\frac35.
\tag{CS7}
\]

K převodu této jediné meze na hybnosti stačí další dvě přesné Taylorovy
nerovnosti:

\[
\begin{aligned}
\cos\frac{93}{100}
&\le1-\frac12\left(\frac{93}{100}\right)^2
+\frac1{24}\left(\frac{93}{100}\right)^4\\
&=\frac{1436925201}{2400000000}<\frac35,
\end{aligned}
\tag{CS8}
\]

\[
\begin{aligned}
\cos\frac9{20}
&\ge1-\frac12\left(\frac9{20}\right)^2
+\frac1{24}\left(\frac9{20}\right)^4
-\frac1{720}\left(\frac9{20}\right)^6\\
&=\frac{4610288951}{5120000000}>\frac9{10}.
\end{aligned}
\tag{CS9}
\]

Z (CS7),(CS8) plyne \(p_c<\arccos(3/5)<93/100\). Z (CS2),(CS7),(CS9)
plyne \(\cos\kappa_c<9/10<\cos(9/20)\), tedy \(\kappa_c>9/20\). Spolu s
již doloženou kladností \(K_c\) dostáváme

\[
\boxed{0<K_c<\frac{93}{100}-\frac9{20}=\frac{12}{25}.}
\tag{CS10}
\]

Nepoužili jsme nové přibližné řešení kontinuální rovnice ani nový odhad
jejích derivací.

\hypertarget{rozestup-na-celuxe9m-kruhu-a-dux16fsledek-pro-uxfaplnou-periodu}{%
\subsection{Rozestup na celém kruhu a důsledek pro úplnou
periodu}\label{rozestup-na-celuxe9m-kruhu-a-dux16fsledek-pro-uxfaplnou-periodu}}

Pro stejné znaménko křivek dávají (CS5),(CS10)

\[
\boxed{K_b-K_c>\frac{29}{60}-\frac{12}{25}=\frac1{300}.}
\tag{CS11}
\]

Oba kladní reprezentanti leží v \((0,\pi)\), takže jejich rozdíl je již
kruhovou vzdáleností. U opačných znamének je kruhová vzdálenost
\(\min(K_b+K_c,\,2\pi-K_b-K_c)\). První člen je větší než \(29/60\);
druhý je díky \(K_b<\pi\) a (CS10) větší než \(\pi-12/25>1/300\).
Zrcadlové dvojice proto nevytvářejí žádnou shodu přes šev kruhu.

Pro všechna \(e\in I_e\) a všechna znaménka \(s,t\in\{-1,1\}\) tedy
platí

\[
\boxed{
\operatorname{dist}_{\mathbb R/(2\pi\mathbb Z)}
\bigl(sK_b(e),tK_c(e)\bigr)>\frac1{300}.
}
\tag{CS12}
\]

Tím je vyloučena současnost těchto dvou tříd fyzikálních stacionárních
singularit úplného \(F_2\) při stejném \((K,E)\). Jde o vzájemné
oddělení již klasifikovaných vázaných a běžných kontinuálních křivek;
samotné lemma nepřiděluje číselné poloměry všem komplexním obvodům ani
novou globální Fourierovu konstantu. Původní vstup, vybraný jev a
normalizace zůstávají zachovány.

\section{Společný pólový a kontinuální příspěvek druhého kontaktu}

\textbf{NON-CANONICAL / L1. Analytický pracovní výsledek.}

Tento oddíl pojednává o všech osmi fyzikálních křivkách setkání pólu
s hranou, určených v~(PT23)--(PT27). Používá původní úplný sudý
resolvent, včetně celého kontinua a všech kontaktně viditelných pólů.

\subsection{Výsledek a jeho přesný rozsah}

V každém z osmi přípustných setkání pólu s hranou je
\textbf{společný místní pólový a kontinuální příspěvek k vnějšímu
\(F_2\) holomorfní v~\((K,E)\)}. Singularita integrandu se při
setkání stává odstranitelnou v diferenciálu na rozvinuté ploše.
Blízký pól a bod větvení mají stejný vnější posun, takže je lze
obejít jediným společným fyzikálním obvodem. Žádné pólové reziduum
se samostatně nevynechává ani nepřičítá.

Tvrzení odstraňuje tuto konkrétní třídu zdánlivých singularit.
Samo však neříká, že celé \(F_2\) je holomorfní v každém bodě
těchto křivek: při stejném \((K,E)\) může v jiném \(P\) ležet
další stacionární vázaný pól nebo jiný kritický bod. Přesněji je
\(F_2\) součtem zde dokázaného holomorfního místního příspěvku
a zbývající periody. Nemá-li zbývající perioda současné fyzikální
sevření obvodu, je tam holomorfní i celé \(F_2\). Místní
transverzalita nenahrazuje globální kontrolu ostatních stacionárních
vázaných pólů.

Níže uvedená místní věta je úplná. Závěrečná podsekce přesně
vymezuje další kvantitativní údaj potřebný pro užitečný číselný
globální odhad regularity: doloženou společnou velikost okolí
místních souřadnic a obvodů, následovanou kontrolou ostatních
fyzikálních kritických bodů. Nevyžaduje se nedokázaná hladkost
samostatně odděleného pólového nebo kontinuálního členu.

\subsection{Původní vzorce a označení}

Pišme \(e=E-\beta\), \(k=K-P\) a zvolme pulzní znaménko
\(\sigma\). V okolí setkání jsou pulzní vlastní hodnoty různé,
takže zvolená pulzní fáze i váha jsou holomorfní:
\[
\omega=e-\sigma\varepsilon(K-P),\qquad
h_\sigma=\frac{1-\sigma\varepsilon'(K-P)}2.
\]
Položme
\[
v=\cos P,\qquad x=e^{-i(\omega-g)},\qquad t=v-5\cos\omega,
\]
\[
C=x^2-\frac25vx-\frac35,
\qquad B=\frac35x^2+\frac25vx-1,
\]
\[
D_o=xC+qB,\qquad
N=x^2+\frac35(v+q)x+vq,\qquad
f_A=\frac{2xN}{D_o}-1.
\tag{PE1}
\]
Objektový kořen splňuje \(q+q^{-1}=t/2\). Místní diferenciál
druhého kontaktu je přesně
\[
\Omega=\frac1{2\pi}h_\sigma f_A\,\mathrm dP.
\tag{PE2}
\]

Znaménko \(s=+1\) nebo \(-1\) označuje objektovou hranu
\(t=4s\); je odlišné od pulzního znaménka \(\sigma\). Zaveďme
\[
c_0=\cos(g/2),\qquad s_0=\sin(g/2),\qquad
\xi=e^{ig/2},\qquad v_a=5c_0-4.
\]
Dvě objektová setkání jsou
\[
(s,v_0,\omega_0)=(-1,v_a,g/2),
\qquad (+1,-v_a,\pi+g/2).
\tag{PE3}
\]
Pro každé znaménko \(P\) a \(k\) platí energetická rovnice
\(e_0=\omega_0+\sigma\varepsilon(k_0)\) a
\(K_0=P_0+k_0\), přičemž \(\sigma=+1\) pro první skupinu
a \(\sigma=-1\) pro druhou. Jde přesně o osm křivek již určených
v~(PT25).

\subsection{Jednotně nenulové derivace ve všech osmi setkáních}

Označme \(v_{\rm edge}=\omega_{\rm edge}'(P_0)\) a
\(v_{\rm pulse}=\sigma\varepsilon'(k_0)\). Dokázané fyzikální
meze jsou
\[
|v_{\rm edge}|<\frac38,\qquad
|v_{\rm edge}-v_{\rm pulse}|>\frac9{40},\qquad
|v_a|<\frac18,\qquad s_0>\frac12.
\tag{PE4}
\]
Přesný pulzní vzorec rovněž dává
\(|\varepsilon'|\le2/\sqrt5<9/10\). Při pevných \(K,e\)
je úplná derivace podle \(P\)
\[
t_P=-\sin P+5\sin\omega\,\sigma\varepsilon'(k)
=5\sin\omega\,(v_{\rm pulse}-v_{\rm edge}).
\]
Proto na všech osmi křivkách setkání jednotně platí
\[
\boxed{\frac9{16}<|t_P|<6.}
\tag{PE5}
\]
Ostrá horní mez plyne již z~\(1+5(9/10)<6\). To je potřebná
prostorová transverzalita pro místní souřadnici na rozvinuté ploše.
Nejde o tvrzení, že objektový pól v mezeře má mimo toto místní
setkání nenulovou energetickou derivaci.

V bodech (PE3) platí
\[
x_0=-s\xi,\qquad q_0=s.
\]
Přímé dosazení do původních polynomů dává
\[
\boxed{
B_0=\frac85\xi(\xi-1),\qquad
N_0=(5c_0-3)(\xi-1).
}
\tag{PE6}
\]
Zvláště
\[
|B_0|=\frac{16}{5}\sin(g/4)
=\frac45(\sqrt5-1)>\frac45,
\qquad N_0\ne0.
\tag{PE7}
\]
Tyto přesné identity vylučují pól vyššího řádu v uniformizační
souřadnici a ukazují, že objektový práh je skutečnou viditelnou
rezonancí. Nenahrazují celé \(f_A\) samotným pólem.

\subsection{Přesná místní uniformizace a pohybující se pól}

Zvolme jedno setkání a nechme \(\zeta=(K,e)\) měnit v malém
komplexním okolí \(\zeta_0=(K_0,e_0)\). Použijme
uniformizační souřadnici \(\theta\), definovanou vztahy
\[
q=s e^{-i\theta},\qquad t=4s\cos\theta.
\tag{PE8}
\]
Rovnice
\[
t(P,K,e)-4s\cos\theta=0
\]
má podle (PE5) jediné holomorfní řešení
\(P=p(\theta,\zeta)\). Toto řešení je sudé v~\(\theta\)
a reálné pro reálné \(\theta,\zeta\). Definujme bod větvení
\(P_e(\zeta)=p(0,\zeta)\). Derivováním dostaneme přesně
\[
p_\theta=-\frac{4s\sin\theta}
{t_P(p(\theta,\zeta),\zeta)}
=\theta L(\theta,\zeta),\qquad
L(0,\zeta_0)=-\frac{4s}{t_P(P_0,\zeta_0)}\ne0.
\tag{PE9}
\]

Vyjádřeme \(D_o\) v těchto souřadnicích, včetně původních
\(x,v\) z~(PE1). V bodě setkání je \(p_\theta=0\), a tedy
\[
\partial_\theta D_o(0,\zeta_0)=-isB_0\ne0.
\tag{PE10}
\]
Holomorfní věta o implicitní funkci a dělení jednoduchou nulou
proto dávají
\[
\boxed{
D_o(\theta,\zeta)=d(\theta,\zeta)
[\theta-a(\zeta)],\qquad
d\ne0,\qquad a(\zeta_0)=0.
}
\tag{PE11}
\]
V tomto dostatečně malém souřadnicovém okolí na rozvinuté ploše
je právě jedna nula \(D_o\). Její místní průmět do proměnné
\(P\) je \(P_b(\zeta)=p(a(\zeta),\zeta)\). Sudost
\(p\) dává silnější identitu
\[
\boxed{
P_b(\zeta)-P_e(\zeta)
=a(\zeta)^2J(a(\zeta)^2,\zeta),\qquad
J(0,\zeta_0)=-\frac{2s}{t_P(P_0,\zeta_0)}\ne0.
}
\tag{PE12}
\]
Pól a bod větvení se tedy spojují kvadraticky v příčné pólové
souřadnici. Zacházení s nimi jako se dvěma nezávislými kořeny
s opačnými posuny by tuto strukturu přehlédlo.

\paragraph{Reálnost a fyzikální strana vznikajícího pólu.}

Pro reálné \(\theta,\zeta\) položme
\(\eta=\eta_0+\theta\), kde \(\eta_0=\pi\) pro
\(s=-1\) a \(\eta_0=0\) pro \(s=+1\). Původní vzorec
\[
\begin{aligned}
D_o={}&\frac{4i}{5}e^{-i(\eta+3(\omega-g))/2}\,
\bigg[
(\cos(\omega-g)-v)
\sin\frac{\eta-(\omega-g)}2\\
&\hspace{43mm}
-4\sin(\omega-g)\cos\frac{\eta-(\omega-g)}2
\bigg]
\end{aligned}
\]
je nenulový analytický fázový činitel krát reálná analytická
funkce. Její jednoduchá nula v~\(\theta\) zaručuje, že
\(a(\zeta)\), a tedy i \(P_b(\zeta)\), jsou reálné pro
reálné \(\zeta\). Nejde o předpoklad, že libovolné komplexní
póly zůstávají reálné.

V přísné místní mezeře má fyzikální kořen určený vnějším
jednostranným pokračováním znaménko
\(\operatorname{sgn}(s\sin\theta)=\operatorname{sgn}(\sin\omega)\).
V obou setkáních je
\(\operatorname{sgn}(\sin\omega_0)=-s\). Fyzikální pól
v mezeře proto odpovídá přesně podmínce
\[
\boxed{a(\zeta)<0.}
\tag{PE13}
\]
Pokračování s~\(a>0\) leží na druhém objektovém listu.
Při \(a=0\) je objektová singularita odmocninového typu
a nemá atom.

\subsection{Společný normální tvar úplného diferenciálu}

Vyjádřeme úplný diferenciál (PE2) v souřadnici \(\theta\).
Definujme
\[
A(\theta,\zeta)
=\frac{h_\sigma\,2xN\,L}{2\pi d}.
\]
Všechny činitele jsou v souřadnicovém okolí holomorfní.
Podle (PE9)--(PE11) platí
\[
\Omega=
\left[
A(\theta,\zeta)\frac{\theta}{\theta-a(\zeta)}
-\frac{h_\sigma\theta L}{2\pi}
\right]\mathrm d\theta.
\]
Analytická dělená diference
\[
\frac{A(\theta,\zeta)-A(a(\zeta),\zeta)}
{\theta-a(\zeta)}
\]
ukazuje, že lze tento diferenciál přesně napsat ve tvaru
\[
\boxed{
\Omega=H(\theta,\zeta)\,\mathrm d\theta
+\frac{a(\zeta)R(\zeta)}{\theta-a(\zeta)}\,\mathrm d\theta,
}
\tag{PE14}
\]
kde \(H\) je holomorfní a
\(R(\zeta)=A(a(\zeta),\zeta)\) je holomorfní. Při setkání
reziduum zaniká a diferenciál je holomorfní i přes
\(\theta=0\). Jde o odstranitelnost \textbf{diferenciálu na
rozvinuté ploše}, nikoli o tvrzení, že samotné \(f_A\) je
omezené v~\(P\).

Vedoucí koeficient je explicitní. Z~(PE6) vychází
\[
\boxed{
R(\zeta_0)=
\frac{i h_\sigma(P_0,K_0)\,5s(5c_0-3)}
{2\pi t_P(P_0,K_0,e_0)}.
}
\tag{PE15}
\]
Nerovnosti \(h_\sigma>1/20\), \(5c_0-3>1\), (PE5)
a \(3<\pi<10/3\) dávají
\[
\boxed{\frac1{160}<|R(\zeta_0)|<3.}
\tag{PE16}
\]
Reziduum tedy v příčné pólové souřadnici zaniká přesně prvního
řádu, nikoli pouze blíže neurčeného řádu.

\subsection{Nenulová příčná souřadnice a nepřítomnost fyzikálního sevření obvodu}

Na samotné hraně se (PE1) rozkládá přesně jako
\[
(D_o)_{\rm edge}
=2x(x-s)[\cos(\omega-g)-\cos\omega].
\]
Derivování tohoto výrazu v~(PE3) s použitím (PE10) dává
\[
\left.
\frac{\partial a}{\partial\omega_{\rm edge}}
\right|_{\zeta_0}
=\frac52(1+c_0).
\tag{PE17}
\]
Položme \(\Delta v=v_{\rm edge}-v_{\rm pulse}\).
Derivování hranové rovnice dává
\[
\partial_eP_e=\frac1{\Delta v},\qquad
\partial_e\omega_{\rm edge}=\frac{v_{\rm edge}}{\Delta v}.
\tag{PE18}
\]
Původní mez \(|v_a|<1/8\) implikuje \(|\sin P_0|>7/8\),
a tedy \(|v_{\rm edge}|>7/40\). Společně s~(PE4) tyto
vzorce dokazují
\[
\boxed{
\frac8{11}<|\partial_eP_e(\zeta_0)|<\frac{40}{9},
\qquad
\frac12<|\partial_ea(\zeta_0)|<\frac{25}{3}.
}
\tag{PE19}
\]
U druhé dolní meze lze konkrétně použít
\((9/2)(7/55)=63/110>1/2\). Zvláště je \(a\) platnou
místní souřadnicí příčnou ke každé rezonanční křivce.

Rovnice (PE12) nyní dává rozhodující identitu
\[
\boxed{
\partial_eP_b(\zeta_0)=\partial_eP_e(\zeta_0)\ne0.
}
\tag{PE20}
\]
Obě funkce jsou holomorfní a pro reálné parametry reálné.
Po zmenšení reálného parametrického okolí mají jejich derivace
podle \(e\) stejné znaménko a absolutní hodnotu větší než
\(1/2\). Při původním vnějším předpisu
\(e\mapsto e+i\epsilon\), \(\epsilon>0\), proto oba
průměty singularit leží pro dostatečně malé \(\epsilon\)
ve stejné polorovině proměnné \(P\):
\[
\operatorname{sgn}\operatorname{Im}P_b(K,e+i\epsilon)
=\operatorname{sgn}\operatorname{Im}P_e(K,e+i\epsilon)
=\operatorname{sgn}(\Delta v).
\tag{PE21}
\]
Znaménkové tvrzení plyne přímo integrací reálné části holomorfní
derivace podle \(e\) podél krátké svislé úsečky. Nepředpokládá,
že lze na jedné straně setkání pólový člen zahodit.

\subsection{Pevný fyzikální obvod a holomorfní místní perioda}

Zvolme dostatečně malý pevný reálný interval
\([P_0-r,P_0+r]\). Jeho konce míjejí bod větvení i pól.
Místní souřadnicové okolí neobsahuje jiný objektový bod větvení
ani jinou nulu \(D_o\); pulzní větve jsou v něm různé.
Zmenšeme parametrické okolí tak, aby \(P_e,P_b\) zůstaly
uvnitř disku \(|P-P_0|<r/4\).

Nahraďme interval pevným půlkruhovým obloukem poloměru \(r\)
v polorovině opačné ke společnému znaménku (PE21). Pro všechna
blízká reálná \(K,e\) a malé kladné vnější \(\epsilon\)
neprotne deformace reálného intervalu na tento oblouk ani jeden
průmět singularity. Podél deformace pokračujme původním objektovým
listem. Jde o jedinou homotopii úplného integrandu; neprovádí se
samostatná volba znaménka odmocniny ani samostatné odečtení
rezidua vázaného pólu.

Tentýž pevný oblouk, zvednutý na původní list, míjí všechny místní
singularity i pro komplexní parametry v dostatečně malém
polydisku. Jeho integrand je holomorfní a na kompaktní cestě
jednotně omezený. Perioda je proto holomorfní v~\(K\) i~\(e\).
Jednostranný limit \(\epsilon\downarrow0\) ji ztotožňuje
s původním vnějším místním příspěvkem, včetně fyzikálního
pólového příspěvku při \(a<0\).

Normální tvar (PE14) činí tuto analytičnost explicitní.
Nechť \(\theta_-(\zeta),\theta_+(\zeta)\) jsou zvednuté
hodnoty konců a \(F\) je holomorfní primitivní funkce k~\(H\)
v proměnné \(\theta\). Oba konce zůstávají oddělené od nuly.
S větvemi logaritmu a celým číslem oběhu \(n\), které jsou
určeny původní vnější homotopií, je místní perioda
\[
\begin{aligned}
I_{\rm loc}(\zeta)={}&
F(\theta_+(\zeta),\zeta)-F(\theta_-(\zeta),\zeta)\\
&+a(\zeta)R(\zeta)
\big[\log(\theta_+(\zeta)-a(\zeta))
-\log(\theta_-(\zeta)-a(\zeta))+2\pi i n\big].
\end{aligned}
\tag{PE22}
\]
Argumenty logaritmů jsou v základním bodě nenulové a zůstávají
takové i v okolí. Každý člen (PE22) je holomorfní přes
\(a=0\). Nevzniká singularita typu \(a\log a\). Koeficient
možné opravy počtem oběhů je přesně zanikající reziduum \(aR\);
celé číslo je určeno fyzikálním jednostranným pokračováním,
nevolí se za účelem zlepšení regularity nebo pravděpodobnostní
meze. Tím je dokázána společná místní věta uvedená na začátku
oddílu.

\subsection{Vztah k úplnému druhému kontaktu a soupisu pólů}

Rozdělme přesnou vnější periodu na tento místní příspěvek
a zbývající intervaly \(P\) a pulzní větev. Pokud při stejných
parametrech nastane několik z osmi setkání, lze každý místní
příspěvek ošetřit ve vlastních souřadnicích. Různé pulzní větve
zůstávají oddělenými analytickými sčítanci, protože tyto hodnoty
\(k\) nejsou pulzní degenerací. Výsledkem zůstává původní úplný
druhý kontakt.

Jedinou další podmínkou pro holomorfnost celého \(F_2\) ve
zvolených parametrech je existence pravidelného fyzikálního
obvodu pro zbývající periodu. Jde o podmínku týkající se jiných
možných stacionárních bodů, nikoli o zbylou neurčitost společného
místního normálního tvaru pólu a hrany. Stacionární póly společné vázané a pulzní disperze jsou nyní
úplně klasifikovány v (BS25). Pro globální odhad je nutné
spojit jejich místní příspěvky s ostatními částmi původní periody.

Pro samotný objektový pól tatáž souřadnice ověřuje místní stranu
vzniku větve z nezávislého soupisu pólů. Podél objektové hrany je
\(\mathrm d\omega_{\rm edge}/\mathrm dv=-1/(5\sin\omega_0)\),
a proto (PE17) dává
\[
\left.\frac{\mathrm da}{\mathrm dv}\right|_0
=-\frac{1+c_0}{2\sin\omega_0}.
\tag{PE23}
\]
Společně s~\(a<0\) pro fyzikální kořen to znamená, že první
nízká větev vzniká na \(v>v_a\), zatímco její opačná kopie
na \(v<-v_a\). Samo toto místní tvrzení nedokazuje globální
počet objektových vázaných větví.

\subsection{Co zbývá k explicitním místním konstantám}

Všechny podmínky použité místní větou jsou ve skutečných
setkáních ověřeny, s jednotnými nenulovými dolními mezemi
(PE5), (PE7), (PE19). Analytičnost a kompaktnost osmi křivek
nad uzavřeným energetickým oknem proto dávají společné kladné
souřadnicové a parametrické poloměry.
\textbf{Číselné hodnoty těchto společných poloměrů zde nejsou
doloženy.}

Pro číselné meze derivací by postačoval konkrétní poloměr
oblouku \(r\) v~\(P\) a komplexní parametrický poloměr
\(\rho\), pro které explicitní místní souřadnice na všech
osmi rodinách splňují
\[
\frac16\le|J|\le8,\qquad |d|\ge\frac25,\qquad
|P_e-P_0|\le\frac r4,\qquad
|a|\le\frac12\sqrt{\frac{3r}{32}},
\tag{PE24}
\]
a zvednutý oblouk zůstává v souřadnicovém okolí. Dále postačují
elementární meze \(|h_\sigma|\le2\), \(|x|\le2\),
\(|N|\le13\). Jde o konečné nerovnosti pro původní explicitní
analytické funkce, nikoli o nový předpoklad spektrální hladkosti.
Jejich ostré verze platí v nějakém okolí podle dokázaných mezí
v základním bodě.

Na tomto oblouku je \(|P-P_e|\ge3r/4\). Rovnice (PE12),
použitá jako místní sudý projekční vzorec
\(P-P_e=\theta^2J(\theta^2,\zeta)\), implikuje
\[
|\theta|\ge\sqrt{\frac{3r}{32}},\qquad
|\theta-a|\ge\frac12\sqrt{\frac{3r}{32}},\qquad
|D_o|\ge\frac15\sqrt{\frac{3r}{32}}.
\]
Délka oblouku je \(\pi r\). Úplná místní perioda proto má
explicitní podmíněnou mez
\[
|I_{\rm loc}|
\le r\left[260\sqrt{\frac{32}{3r}}+1\right]
<1041\sqrt r\qquad(0<r\le1).
\tag{PE25}
\]
Cauchyův vzorec následně dává například
\[
|\partial_E^j I_{\rm loc}|
\le j!\,1041\sqrt r\,\rho^{-j}
\tag{PE26}
\]
ve středu příslušného parametrického disku. Pro vyhodnocení
uvnitř disku se za \(\rho\) vezme skutečná dostupná vzdálenost
od jeho hranice. Pro tento místní člen tedy zbývá vyčíslit společnou velikost
okolí splňujícího (PE24).
Globální zlepšení časového odhadu navíc vyžaduje kontrolu
ostatních fyzikálních kritických bodů popsanou výše.

Již dokončený důkaz konečného času na tomto místním zpřesnění
nezávisí. Výsledek zachovává původní zdroj, pevný počátek času,
vybraný jev, normalizaci, štítky i paměť a referenci. Nedokládá
nově realizovaný konečný laboratorní experiment ani vnější
nezávislé přijetí tohoto dodatečného analytického lemmatu.


\section{Přesný stav důkazu a další číselný krok}

\subsection{Co tento dodatek uzavírá}

Fázová rovnice určuje všechny fyzikální kontaktně viditelné
objektové póly na celém kruhu. U hlavní kladné a obou nízkých
větví je v původním okně vyloučena shoda rychlostí s pulzem.
Zbývající hlavní větev posunutá o \(\pi\) dává přesně dvě
konfigurace (BS25). Pro obě platí
\[
\partial_P^2[\omega_{\rm lower}(P)-\varepsilon(K-P)]
>\frac{19}{64}>\frac14,\qquad
|t|<2,\qquad r_b>\frac1{32},\qquad h_\sigma>\frac38.
\]
Jejich původní objektová hybnost má \(|P|>\pi/6>1/100\).
Póly proto leží mimo malé vybrané výstupní okno, avšak zůstávají
součástí úplného \(F_2\) a úplné časové chyby. Rozestup od obou
běžných kontinuálních stacionárních křivek je větší než \(1/300\).

Ve všech osmi setkáních objektového pólu s kontinuální hranou
je společný místní příspěvek holomorfní. Přesná normalizovaná
diferenciální forma zachovává oba členy. Výsledek se nepřevádí
na nepodložené tvrzení o samostatném pólovém členu.

\subsection{Co je ještě třeba vyčíslit pro kratší čas}

Dosavadní obecný diskriminační důkaz z [S4] poskytuje konečný
čas bez následujících zpřesnění. Pro výrazně silnější odhad
založený na skutečných fyzikálních prahech zbývají tři konkrétní
úlohy:
\begin{enumerate}
\item Vybrat společná konkrétní \(r,\rho>0\) pro osm místních
okolí v (PE24), včetně oddělení ostatních kořenů a původního
pulzního listu. Nenulové dolní meze potřebných derivací jsou
již doloženy; samotné číselné poloměry zde stanoveny nejsou.
\item U obou vázaných stacionárních bodů odvodit kvantitativní
místní kvadratické souřadnice z (BS22), včetně horních mezí
dalších derivací a původních reziduí. Stejně spojit již
klasifikované běžné kontinuální prahy, objektové kuželové
konce a výměnu pulzních větví do úplného vnějšího \(F_2\).
Rozestup (CS12) vylučuje souběh vázaného a běžného
kontinuálního stacionárního bodu.
\item Přenést společné místní odhady na celý sloupec
\(V=(-2F_2,4)^T/(4+F_1F_2)\), započítat ostré energetické
konce v původní normě a vyčíslit součet časového chvostu.
Odhad samotného objektového pólu nebo samotné reálné hustoty
k tomuto kroku nestačí.
\end{enumerate}

Jde o konkrétní kvantitativní úlohy pro již dané funkce.
Tento dodatek neprokazuje, že je třeba nový fyzikální nebo
spektrální předpoklad. Současně žádnou chybějící konstantu
nenahrazuje pouhou existencí konvergence.

\subsection{Ostré energetické konce mají přesně placenou normovou chybu}

Pro další postup lze použít pomocný hladký skalár
\(0\le q_\eta(E)\le1\), který je roven jedné na
\([E_*-h+\eta,E_*+h-\eta]\), je podepřen uvnitř původního
\(I\) a plynule klesá k nule v obou krajních proužcích.
Vektor se nepřenormovává. Přesná rovnoměrná energetická míra
původního zdroje dává
\[
\nu^2=\|(1-q_\eta(H_0))\psi\|^2
=\frac1{2h}\int_I|1-q_\eta(E)|^2\,\mathrm dE
\le\frac{\eta}{h}.
\]
Při návratu od pomocného vektoru k původnímu vstupu přidá
porovnání dvou izometrických vývojů nejvýše
\[
2\nu\le2\sqrt{\eta/h}.
\]
Například volba \(\eta=10^{-18}\), při \(h=10^{-3}\),
má vrácenou normovou chybu
\[
2\nu\le2\sqrt{10}\,10^{-8}<6.4\cdot10^{-8}.
\]
Tato volba sama neurčuje čas, protože derivace přechodu
rostou s \(\eta^{-1}\). Ukazuje však přesně, jak se ostré
hranice zachovají v konečném účtu. Projektor \(Q\), všechny
štítky i korelace s pamětí a referencí zůstávají původní.

\subsection{Dosavadní konečný čas a pravděpodobnost}

Tento dodatek nemění časy samostatného návrhu [S4]:
\[
N_-=10^{16},\qquad N_+=2^{81\,000\,000\,000\,000}.
\]
Jde o počty původních diskrétních makrokroků. V [S4] je při
těchto časech, za převzatých operátorových a rozptylových
předpokladů, odvozena úplná časová chyba včetně přijímače
\[
\varepsilon_-(N_-)+\varepsilon_+(N_+)
<2.310002\cdot10^{-7}<10^{-6}.
\]
Pro původní \(\varepsilon_S=\varepsilon_D=10^{-6}\)
zůstává celkový amplitudový účet
\[
\left|\sqrt{p_{\rm finite}}-\sqrt{P_0}\right|
\le\varepsilon_S+\varepsilon_D+
\varepsilon_-(N_-)+\varepsilon_+(N_+)
<2.24\cdot10^{-6}.
\]
Původní absolutní dolní mez
\(L=293166031669590845829/2^{96}>(60\cdot10^{-6})^2\)
tedy nadále dává konzervativní
\[
p_{\rm finite}>(57.76\cdot10^{-6})^2
=3.3362176\cdot10^{-9}.
\]
Nevznikl nový menší čas ani nový výpočet absolutní
pravděpodobnosti. Jde o matematický konečný protokol
v původním modelu s původními dodanými operacemi přípravy
a čtení. Neříká se, že byl uskutečněn laboratorní pokus.
Převod makrokroků na sekundy není v modelu doložen.

\section{Veřejná návaznost, kontrola a neměnné zdroje}

Před tímto pokračováním proběhla kontrola pouze čtením
29.~10.~2026 v~10:19:10--10:19:36 CET
(Europe/Rome, UTC+1). PR~\#1817 zůstal otevřeným
nesloučeným návrhem na
\nolinkurl{66ef8220df3ad5764ca941b84223a9f46cf5411f}.
Rezervace \#1816 zůstala otevřená a beze změny.
Nové veřejné pokračování konečného rozptylového času
nebylo nalezeno. Bylo přečteno všech 322 aktuálních hlav;
ověřena je shoda relevantních pinů. Netvrdí se nově
provedené porovnání všech hlav se starým soupisem.

Hlava hlavní větve byla
\nolinkurl{a193e389fb1ddb0bdc89165b015c95a17eb1e425}.
Soubor STATUS nadále označoval Public Canon v106 jako
ACTIVE, s obsahovým commitem
\nolinkurl{9242db4f040d67d9417eacd06612230b74248cff}.
Oddělený provozní záznam \#1808 uváděl, že také druhé
úplné přehrání po vydání bylo zastaveno po 40 minutách.
Není to matematický vyvracející výsledek, ale tato úplná
kontrola po vydání není dokončená. Úspěšné původní
tagové přehrání a kontroly vydaných souborů tím nejsou
zaměňovány s dokončením pozdějšího celého běhu.

Nové algebraické důkazy byly v této pracovní relaci
přezkoumány z několika směrů. Zvláštní kontrola ověřila
parametrizaci fyzikální vrstvy, všechny ostré nerovnosti
v (BS1)--(BS31), počet a znaménka kořenů i kruhový
rozestup (CS12). Jde o vnitřní analytický přezkum.
Výsledek není externím přijetím a nepatří do dřívějšího
auditu 23 skupin a 130 podmínek.

Práce je samostatná, \textbf{NON-CANONICAL / L1}.
Rezervace \#1816 ani veřejný Canon nebyly měněny.
Nebyl proveden veřejný zápis, nový numerický průzkum,
CAS vyhodnocení ani vědecký běh.

\section*{Odkazy na zdroje}

\begin{description}[leftmargin=9mm,style=nextline]
\item[S1] Původní absolutní zachycení, neměnný
commit \nolinkurl{32306016908c0097650277aa416fc0152b2f81a4},
část I, zejména oddíly 1, 4--8 a 11:
\url{https://github.com/mathorn1973/twist-j/blob/32306016908c0097650277aa416fc0152b2f81a4/notes/C-BOUND-PAIR-ABSOLUTE-CAPTURE-N/PROOF.md}
\item[S2] Úplný algebraický audit, neměnný
commit \nolinkurl{66ef8220df3ad5764ca941b84223a9f46cf5411f}:
\url{https://github.com/mathorn1973/twist-j/blob/66ef8220df3ad5764ca941b84223a9f46cf5411f/notes/C-BOUND-PAIR-RESOLVENT-AUDIT-N/PROOF.md}
\item[S3] Původní konečná příprava, čtení a časová mezera:
\url{https://github.com/mathorn1973/twist-j/blob/66ef8220df3ad5764ca941b84223a9f46cf5411f/notes/C-BOUND-PAIR-RESOLVENT-AUDIT-N/FINITE_PROTOCOL.md}
\item[S4] Samostatný předchozí analytický návrh této relace,
\textit{Konečný čas rozptylu původního paketu TWIST-J},
29.~10.~2026, 33 stran, zejména oddíly 12--14:
\url{twist_j_konecny_cas.tex}
\item[S5] Veřejné návrhy a oddělená rezervace:
\url{https://github.com/mathorn1973/twist-j/pull/1817},
\url{https://github.com/mathorn1973/twist-j/pull/1815},
\url{https://github.com/mathorn1973/twist-j/issues/1816}.
\item[S6] Provozní stav vydání a hlavní větev:
\url{https://github.com/mathorn1973/twist-j/issues/1808};
\url{https://github.com/mathorn1973/twist-j/blob/a193e389fb1ddb0bdc89165b015c95a17eb1e425/STATUS.md}.
\end{description}


\end{document}
